% =====================================================================
% FishBot v0.4
% Exact Combinatorial Inference and Opponent Modelling in Six-Player
% Literature
% =====================================================================
\documentclass[11pt,a4paper]{article}

\usepackage[utf8]{inputenc}
\usepackage[T1]{fontenc}
\usepackage[margin=1in]{geometry}
\usepackage{amsmath,amssymb,amsthm}
\usepackage{booktabs}
\usepackage{microtype}
\usepackage[colorlinks=true,linkcolor=blue,citecolor=blue,urlcolor=blue]{hyperref}
\usepackage{natbib}
\usepackage{multirow}
\usepackage{xcolor}

\newtheorem{proposition}{Proposition}
\newtheorem{lemma}{Lemma}
\newtheorem{remark}{Remark}

\newcommand{\halfsuit}{\ensuremath{H}}
\newcommand{\Prob}{\ensuremath{\mathbb{P}}}
\newcommand{\Ex}{\ensuremath{\mathbb{E}}}
\newcommand{\ind}[1]{\ensuremath{\mathbf{1}\!\left[#1\right]}}
\newcommand{\sets}{\ensuremath{\Delta}}

\title{\textbf{FishBot v0.4}\\[0.4em]
\large Exact Combinatorial Inference, Opponent Modelling, and the Limits of
the Uniform Posterior in Six-Player Literature}

\author{}
\date{}

\begin{document}
\maketitle

% =====================================================================
\begin{abstract}
Literature (also called Fish) is a six-player, two-team,
imperfect-information card game in which players interrogate opponents for
specific cards and score by declaring the exact distribution of a six-card
half-suit among their own team. It is a team game with hidden information and
no legal communication channel, so neither the two-player zero-sum machinery
that solved poker nor the fully cooperative machinery developed for Hanabi
applies directly.

The previous version of this engine ended by naming its own bottleneck: its
belief sampler was documented as \emph{not uniform} over the deals consistent
with the public record, so every probability it reported inherited an
unquantified bias. This paper removes that bias, and reports what happened.

We first show that hidden-state inference in Literature reduces to counting
assignments in a degree-constrained bipartite system, which is $\#P$-complete
in general but tiny in the dimension that matters here: real positions have
$25.6$ unlocated cards on average but only $4.0$ \emph{distinct candidate
sets}, so cards are exchangeable and the problem collapses onto an $8 \times 6$
count matrix. A dynamic program over that structure returns the exact partition
function, exact marginals, exact uniform sampling and exact joint probabilities,
validated against exhaustive enumeration. The one part it cannot absorb --- the
disjunctive constraints certifying that an asker held a card of the half-suit
--- we handle by importance sampling with closed-form proposal densities, after
measuring why the two obvious alternatives fail.

The result is that the inference is now provably correct, and \textbf{it changed
nothing}: against the previous champion, with the ask objective held fixed, the
set differential was $-0.22$ per duplicate deal-pair, 95\% CI
$[-0.96, +0.52]$.

To find out why, we built a benchmark the previous version lacked: because the
hidden hands exist in simulation, a posterior can be scored directly on how much
probability it places on where the cards actually are. Over 517 decisions and
13{,}640 uncertain-card predictions, the \emph{exactly uniform} posterior was
\emph{worse} at predicting reality than the biased sampler it replaced
(negative log-likelihood $1.362$ against $1.341$ at a matched budget). The
uniform posterior is a correct theorem about a false hypothesis --- that
players' choices carry no information beyond legality --- and the old sampler's
bias happened to encode a crude opponent model.

Replacing the accident with an explicit one --- a single parameter modelling the
hypothesis that a player asks in a half-suit in proportion to how many cards of
it they hold --- beats both on posterior accuracy ($1.325$) and is worth
$\mathbf{+1.85}$ \textbf{sets per duplicate deal-pair} in play, 95\% CI
$[+1.23, +2.46]$, essentially flat over a threefold range of the parameter. For
scale, the entire objective improvement separating the previous champion from
its own baseline was $1.26$ sets per pair.

The resulting engine beats the previous champion by \textbf{+1.85 sets per
deal-pair} $[+1.54, +2.17]$ over 500 duplicate deals, wins every pair against
the two weakest policies in the ladder, and improves agreement with provably
optimal play in exactly solved positions from $67.6\%$ to $72.5\%$. Nothing else
we tried helped: an eleven-term ask objective, a learned half-suit value, a
rollout-regression objective, and a widened exact endgame solver are all
reported as nulls, and the one screening cell that did resolve failed to
replicate twice. We also build the one search design the variance diagnosis
leaves open --- a lookahead over the posterior rather than over sampled
layouts --- and prove that its natural form is exactly greedy at every depth, so
that only the quota coupling between candidates can make such a search
non-trivial. The coupled form is worth about $+0.15$ sets per
deal-pair over 2400 unselected pairs, 95\% CI $[+0.022, +0.285]$ --- but the two
cells we pre-registered as decisive pool to $[-0.019, +0.324]$ and do not
separate it from zero, so we report a small effect that is probably real and is
not demonstrated. The attempt to settle it produced the more useful finding. Three runs of the
identical configuration returned $+0.570$, $-0.044$ and $+0.386$ sets per
deal-pair, disagreeing by more than their own intervals allow ($Q = 7.03$,
$p = 0.030$), which suggested a between-run variance component that no interval
in this paper models. We measured that directly --- 24 blocks of an A/A, where
the true differential is exactly $0$ --- and it does not exist: coverage of the
nominal 95\% interval is $23/24$, $Q = 18.41$ on 23 df at $p = 0.735$, and the
estimated between-run standard deviation is $0.0000$. The intervals are honest
and the three-cell disagreement was the one-in-thirty event. Establishing that
required giving the harness the ability to observe its own null, which it did not
have: under seat-based seeding an A/A is bit-identical in both halves of a
duplicate pair and its differential is exactly zero for every deal. We also report several
methodological corrections, including a bias in the previous harness's handling
of timed-out games that we demonstrate can place the true value outside a
reported 95\% confidence interval, and a measurement showing that the previous
version's 600-pair experiment cells could not resolve effects below about
$0.44$ sets per pair --- so its smaller null results were uninformative rather
than negative.
\end{abstract}

% =====================================================================
\section{Introduction}
\label{sec:intro}

Literature --- also called Fish or Canadian Fish --- is a six-player,
two-team, imperfect-information card game. Players interrogate opponents for
specific cards and score by \emph{declaring} the exact distribution of a
six-card half-suit among the members of their own team. It sits in a corner of
the game-AI landscape that is unusually poorly served: it is a team game with
hidden information and \emph{no legal communication channel}, so neither the
two-player zero-sum machinery that solved poker nor the fully cooperative
machinery developed for Hanabi applies directly.

This paper is the fourth iteration of an engine for the game. The previous
version, v0.3, established three results that frame everything here. First,
exact belief tracking dominates: an agent that merely chains every logical
implication of the public record beats a public-information heuristic by
$11.86$ sets per duplicate deal-pair. Second, and more surprisingly, search
does not help, and the reason is measurable: the standard deviation of a
position's value across possible hidden layouts is about $2.4\times$ the gap
between the best and the worst candidate move, so any search that evaluates
different candidates against different sampled layouts is ranking luck. Third,
acting on that diagnosis --- by improving the \emph{objective} rather than the
depth of search --- produced the study's largest single gain.

v0.3 closed by naming its own bottlenecks, and this paper attacks them
directly. Its belief sampler was documented as \emph{not uniform} over the set
of deals consistent with the public record, so every probability the engine
reported inherited an unquantified bias. Its ask objective was a hand-weighted
sum of three terms. Its claim policy was a fixed confidence threshold. The
obvious research programme is therefore: make the inference exact, widen and
then learn the objective, and make claiming decision-theoretic.

We carried that programme out, and the headline result is not the one we
expected.

\paragraph{Contributions.}
\begin{enumerate}
\item \textbf{Exact combinatorial inference.} We show that hidden-state
inference in Literature reduces to counting assignments in a
degree-constrained bipartite system, and that although the general problem is
$\#P$-complete, the specific instances arising in this game are tiny in the
dimension that matters: the number of \emph{distinct candidate sets} is on
average $4.0$ even when $25.6$ cards are unlocated. We give an exact dynamic
program over that structure that returns the partition function, exact
marginals, and exact uniform sampling, and validate it against exhaustive
enumeration (Section~\ref{sec:counting}).

\item \textbf{A measured account of the disjunctive constraints.} The
constraint set contains clauses (``the asker held at least one card of that
half-suit'') that the dynamic program does not absorb. We quantify three ways
of handling them and show why the obvious two fail, then give an importance
sampler whose proposal density is computable in closed form and whose
self-normalised weights are consistent for the target
(Section~\ref{sec:clauses}).

\item \textbf{Posterior quality measured against ground truth.} Every belief
metric in the previous version was indirect. Because the hidden hands exist in
the simulator, we can score a posterior directly on how much probability it
places on where the cards actually are. We report negative log-likelihood,
Brier score, top-1 accuracy and calibration for every inference configuration
(Section~\ref{sec:posterior-quality}).

\item \textbf{The central negative-then-positive result.} Making the posterior
\emph{exactly} correct under the standard assumption --- that players' choices
carry no information beyond legality --- made it measurably \emph{worse} at
predicting reality than the biased sampler it replaced, and produced no
improvement in play. The reason is identifiable: the old sampler's bias was
accidentally informative, encoding a crude opponent model. Replacing the
accident with an explicit, single-parameter model of how players choose which
half-suit to ask in recovers the loss and then beats both
(Sections~\ref{sec:oppmodel} and~\ref{sec:results}).

\item \textbf{A wider, ablatable ask objective}, including two terms that the
literature suggests are unstudied: the \emph{exposure} of one's own already-
public cards to the opponent who would receive the turn on a failed ask, and
the \emph{signalling} value of an ask, since the choice of half-suit is the
only legal communication channel in the game and is simultaneously read by the
opposition (Section~\ref{sec:objective}).

\item \textbf{Methodological corrections} to the previous evaluation harness,
including a bias in how timed-out games were handled that we demonstrate can
put the true value outside a reported 95\% confidence interval
(Section~\ref{sec:eval}); a measurement showing the previous version's
experiment cells could not resolve effects below about $0.44$ sets per pair, so
several of its reported nulls were uninformative rather than negative; and a
worked example of a screening study manufacturing a false positive, caught by
replication (Section~\ref{sec:res-replication}).

\item \textbf{A widened exact solver}, which is the source of two structural
facts about the game: the number of live cards is always exactly six times the
number of unresolved half-suits, and the perfect-information value has the
closed form $\mathrm{sign}(\text{mover}) \times (2f - m)$. The second means the
previous version's headline absolute result --- 100\% agreement with optimal
play in information-resolved positions --- was a shallower certificate than it
read, and we reproduce it on a corpus that is 79\% two-half-suit
(Section~\ref{sec:exact}).

\item \textbf{The search the variance diagnosis leaves open, and why its obvious
form cannot work.} Every search this project has tried samples layouts, and the
measured spread across them is why they all failed. A lookahead over the
posterior instead is deterministic given the belief, so that variance is
structurally absent. We show (Proposition~\ref{prop:greedy}) that the natural
version of it --- one whose transition edits only the asked card's row --- is
\emph{exactly} greedy at every depth by an exchange argument, so it is a no-op by
construction rather than by measurement, and only the quota coupling between
candidates can make such a search non-trivial. Whether the coupled version helps is still
open, and the attempt to resolve it produced a methodological result we consider
more valuable than the search would have been: a direct measurement, over 24
blocks of an A/A, that this harness's duplicate-deal intervals cover at their
nominal rate and carry no between-run variance component --- which required first
making the harness able to observe its own null at all
(Sections~\ref{sec:lookahead} and~\ref{sec:res-lookahead}).
\end{enumerate}

% =====================================================================
\section{Related work}
\label{sec:related}

\paragraph{Literature itself.} We could find no peer-reviewed AI work on
Literature, Canadian Fish or Russian Fish. Several open-source implementations
with bots exist; the best documented uses Q-learning on a four-player variant
with a first-order knowledge encoding. The game family did recently enter a
standardised testbed: \citet{goadrich2026valet} include Go Fish, and report that
``Go Fish and Crazy Eights show little separation between MCTS and random play,
likely because CardStock's determinizations account only for visibility and not
for information inferred from action history.'' That is a published negative
result about exactly the failure mode this paper attacks from the other side ---
they identify that action history carries information their determinizations
ignore; we quantify how much, and put it in.

\paragraph{Determinization and its pathologies.} Evaluating an
imperfect-information position by sampling consistent worlds and solving each
one perfectly is the standard technique, and its two failure modes are named:
\emph{strategy fusion}, where the searcher implicitly assumes it can act
differently in states it cannot distinguish, and \emph{non-locality}
\citep{frank1998search}. \citet{long2010understanding} give the conditions under
which it nonetheless works, parameterised by leaf correlation, bias, and a
disambiguation factor. Two consequences shape this work. First, our
disambiguation factor is unusually high --- every ask is a hard constraint on the
deal --- which is favourable. Second, and decisively, the \textsc{Claim} action
is a strategy-fusion trap: inside a determinization the searcher knows the true
layout, so every claim evaluates as a certain gain and the declared assignment
differs between sampled worlds. Claims in v0.4 are therefore never selected by
search. \citet{cowling2012ismcts} introduce Information Set MCTS and state
plainly that they do not consider belief distributions; v0.3's own experiments
found ISMCTS losing badly to its own prior, which is consistent with both that
and with the variance diagnosis it measured.

\paragraph{Inference from what opponents chose.} The closest published line is
\citet{solinas2019improving} and \citet{rebstock2019policy}, who infer hidden
cards in trick-taking games from the actions players took, using a learned
policy as the likelihood. Their factorised belief does not enforce hand-size
constraints; they note the constraint could be added and do not add it. Our
setting lets us do both at once: the quota constraints are handled exactly by
the dynamic program of Section~\ref{sec:counting}, and the choice likelihood
enters as an importance weight on top. \citet{solinas2023history} prove that
exhaustive history filtering is polynomial only when the public tree is sparse;
Literature's is dense, so their route does not port --- but the count abstraction
makes the constraint side tractable anyway.

\paragraph{Counting and sampling with fixed margins.} The combinatorial problem
underneath is counting bipartite degree-constrained subgraphs, equivalently a
restricted permanent, $\#P$-complete in general \citep{valiant1979complexity},
with a celebrated approximation scheme \citep{jerrum2004polynomial} and a
literature on sequential importance sampling for contingency tables
\citep{chen2005sequential} --- along with cautionary results showing SIS can be
exponentially bad on some instances \citep{bezakova2012negative}. We need none
of that machinery, because the instances here are small in the dimension that
matters. The constraint-programming counterpart is R\'egin's global cardinality
constraint \citep{regin1996generalized}, which achieves arc consistency by a
flow argument; our dynamic program is stronger in that it also \emph{counts},
which is what claim probabilities require.

\paragraph{Team games without communication.} Literature is a two-team zero-sum
game with pre-play agreement and no in-play communication, which is exactly
\citeauthor{celli2018computational}'s regime; the solution concept is a team
maximin equilibrium with correlation, not a Nash equilibrium
\citep{vonstengel1997team,celli2018computational}. The convergence guarantees
that underwrite two-player zero-sum solvers do not transfer. Hanabi
\citep{bard2020hanabi} is the closest studied game in spirit, and one of its
methodological contributions transfers cleanly: cross-play between
independently trained partners as the instrument for detecting conventions that
work only against copies of oneself \citep{hu2020otherplay}. What does not
transfer is the remedy, because Hanabi is fully cooperative while Literature's
only communication channel --- the choice of which half-suit to ask in --- is
\emph{costly} and read by the opposition as well as by one's partner. We found
no published treatment of that combination, and it is the motivation for the
\texttt{signal} term of Section~\ref{sec:objective}.

% =====================================================================
\section{The game}
\label{sec:game}

Six players sit in a circle in alternating team order: seats $\{0,2,4\}$
against $\{1,3,5\}$. The deck is partitioned into \emph{half-suits} of six
cards each. Our primary variant uses 54 cards and nine half-suits: the eight
natural half-suits ($2$--$7$ and $9$--$A$ of each suit) plus a ninth composed
of the four $8$s and two individually identified jokers, a Red Joker and a
Black Joker. Each player is dealt nine cards. The classic 48-card,
eight-half-suit variant is also supported.

On turn, a player either \textbf{asks} or \textbf{claims}.

An \textbf{ask} names an opponent and a specific card. It is legal only if the
asker holds at least one card of that card's half-suit and does not hold the
card itself. If the opponent has it, the card transfers \emph{face up} and the
asker keeps the turn; otherwise the turn passes to the opponent who was asked.

A \textbf{claim} names a half-suit and declares which teammate holds each of
its six cards. If the declaration is exactly right, the claiming team scores
the half-suit. If any of the six is with an opponent, the opposing team scores
it. If all six are within the claiming team but the declared distribution is
wrong, the half-suit is nulled and nobody scores. The requirement to name the
exact distribution, not merely to assert possession, is what makes Literature
a game about inference rather than collection.

A player with no cards on move passes the turn to a teammate who has cards. If
an entire team is out of cards, the other team must claim out the remaining
half-suits without asking.

\begin{remark}[Termination is a property of the players, not the rules]
Nothing in the rules forces progress: two opponents can pass a card back and
forth and return to an identical position, so the state graph contains cycles.
Consequently any claim that ``every game ends'' describes an agent's stalling
rule rather than Literature itself. Our agents carry an explicit progress rule
and our harness imposes an action cap; Section~\ref{sec:eval} explains why the
cap must be \emph{scored} rather than discarded.
\end{remark}

\begin{remark}[Solution concept]
Literature is a two-team zero-sum game with pre-play agreement and no in-play
communication. The appropriate solution concept is therefore a team-maximin
equilibrium with correlation (TMECor) in the sense of
\citet{celli2018computational}, not a Nash equilibrium, and the convergence
guarantees available for two-player zero-sum games do not transfer. We compute
no equilibrium and make no equilibrium claim; every strength statement in this
paper is either relative (against a named opponent) or absolute against an
exactly solved subgame.
\end{remark}

% =====================================================================
\section{Information integrity}
\label{sec:integrity}

An imperfect-information study is worthless if the agents can see the hidden
state, and the failure mode is usually subtle rather than flagrant. The
boundary is enforced four ways.

\paragraph{Structural.} The engine's \texttt{GameState} never reaches a
policy. Policies consume only an \texttt{Observation}: their own hand, the
public event log, public hand counts, resolved half-suits, whose turn it is,
and the rules.

\paragraph{Proved by reconstruction.} \texttt{Observation.reconstruct} rebuilds
a seat's entire view from nothing but (rules, that seat's initial hand, the
public event log). Tests assert the engine-provided observation is
\emph{identical} to that reconstruction at every step of full games, for all
six seats. In an audit conducted for this paper the property was re-verified
over $183{,}750$ seat-plies across eight rule configurations with zero
mismatches.

\paragraph{Proved by invariance.} Every quantity v0.4 computes --- posterior
marginals, all eleven ask-objective terms, and the end-to-end policy action ---
is asserted to be bit-identical when the hidden cards of the other five seats
are permuted into a different layout that is equally consistent with the public
record. Candidate permutations are drawn from the acting seat's own belief
sampler and then validated by an \emph{independent} replay of the entire public
history, so the test does not merely assume the sampler is correct.

\paragraph{Separation of training from acting.} Ground truth may be used to
\emph{score} (Section~\ref{sec:posterior-quality} scores posteriors against the
true hands) and to train; no acting policy ever sees it. That line is what
makes the distinction meaningful rather than decorative.

% =====================================================================
\section{Exact combinatorial inference}
\label{sec:counting}

\subsection{The inference problem, stated exactly}

The observation that shapes the whole engine is that \emph{every card movement
in Literature is public}. A successful ask names the card that moved; a claim
reveals all six locations; nothing changes hands unobserved. It follows that
the current location of every card is a deterministic function of the
\emph{initial deal} and the public log, so all hidden-state inference reduces
to constraints on the initial deal. Those constraints are of exactly three
kinds:

\begin{enumerate}
\item \textbf{Candidate sets.} Per card $c$, a set $S_c \subseteq \{0,\dots,5\}$
of players who may have been dealt it. A failed ask excludes the target;
a successful ask pins the card; asking for a card excludes the asker (under
standard no-bluff rules); a claim pins all six.
\item \textbf{Exact counts.} Every player was dealt exactly $q$ cards ($9$ or
$8$); after propagation, player $p$ has a residual quota $q_p$ of the cards
that remain unlocated.
\item \textbf{Disjunctions.} A legal ask certifies that the asker held
\emph{at least one} card of that half-suit at that moment: an OR-constraint
over the cards of the half-suit not yet publicly located.
\end{enumerate}

Write $F$ for the set of \emph{free} cards (those not publicly located and not
pinned by propagation). A \emph{world} is a map $f : F \to \{0,\dots,5\}$ with
\begin{equation}
f(c) \in S_c \quad \forall c \in F,
\qquad
|f^{-1}(p)| = q_p \quad \forall p,
\label{eq:world}
\end{equation}
that additionally satisfies every OR-constraint. Let $\Omega$ denote the set of
such worlds and $\Omega_0 \supseteq \Omega$ the set satisfying only
\eqref{eq:world}.

\begin{proposition}[Uniform posterior under uninformative choice]
\label{prop:uniform}
Suppose the deal is uniform over all $\binom{54}{9,9,9,9,9,9}$ partitions, and
suppose that at every decision point each player's choice of action depends on
their private hand only through which actions are legal. Then the posterior
over worlds given the public history is uniform on $\Omega$.
\end{proposition}

\begin{proof}[Proof sketch]
Let $h$ denote the public history and $w$ a world. The prior is uniform, so
$\Prob(w \mid h) \propto \Prob(h \mid w)$. Factor $\Prob(h \mid w)$ over the
actions in $h$. Under the stated assumption each factor depends on $w$ only
through the legality of the taken action, and card transfers are deterministic
given $w$ and the action. Hence $\Prob(h \mid w)$ takes the same value for
every $w$ that makes all of $h$ legal and reproduces every observed outcome,
which is precisely $\Omega$, and is $0$ otherwise.
\end{proof}

Proposition~\ref{prop:uniform} is the target that v0.3's sampler was
\emph{intended} to hit and, by its own documentation, missed. It is also, as
Section~\ref{sec:posterior-quality} shows, not the target one should want ---
because its hypothesis is false, and interestingly so.

\subsection{Why the counting is tractable}

Counting the elements of $\Omega_0$ is an instance of counting perfect
matchings in a degree-constrained bipartite graph, equivalently a restricted
permanent, which is $\#P$-complete in general \citep{valiant1979complexity}.
The instances arising here, however, are small in the only dimension that
matters. Measured over $901$ decision points from real games:

\begin{center}
\begin{tabular}{lrrrr}
\toprule
quantity & mean & median & p90 & max\\
\midrule
free cards $|F|$                 & 25.6 & 26 & 43 & 45\\
\textbf{distinct candidate sets} & \textbf{3.9} & \textbf{4} & \textbf{6} & \textbf{8}\\
active OR-constraints            & 5.4 & 4 & 12 & 18\\
cards per OR-constraint          & 3.35 & 3 & 5 & 5\\
\bottomrule
\end{tabular}
\end{center}

Cards sharing a candidate set are \emph{exchangeable}. Grouping them collapses
the problem from an assignment over $|F| \le 45$ labelled cards onto a
\emph{group-count matrix} $k \in \mathbb{N}^{G\times 6}$ with $G \le 8$, where
$k_{gp}$ is the number of cards of group $g$ dealt to player $p$.

\subsection{The dynamic program}

Let group $g$ have candidate set $S_g$ and size $n_g$, so $\sum_g n_g = |F|$.
The number of worlds realising a given $k$ is a product of multinomials, hence
\begin{equation}
Z \;=\; |\Omega_0| \;=\;
\sum_{\substack{k \,:\, \sum_p k_{gp} = n_g \ \forall g \\
\sum_g k_{gp} = q_p \ \forall p, \ \ k_{gp}=0 \text{ if } p \notin S_g}}
\;\prod_{g=1}^{G} \frac{n_g!}{\prod_p k_{gp}!}.
\label{eq:Z}
\end{equation}

We evaluate \eqref{eq:Z} by dynamic programming over \emph{players}, carrying
as state the vector $r \in \prod_g \{0,\dots,n_g\}$ of cards still unassigned
in each group. Define $\Phi_0(r) = \ind{r = (n_1,\dots,n_G)}$ and
\begin{equation}
\Phi_{p+1}(r') \;=\;
\sum_{\substack{\kappa \,:\, \sum_{g \in A_p}\kappa_g = q_p\\ r' = r - \kappa,\ \kappa_g \le r_g}}
\Phi_p(r) \prod_{g \in A_p} \binom{r_g}{\kappa_g},
\qquad A_p = \{g : p \in S_g\},
\label{eq:dp}
\end{equation}
so that $Z = \Phi_6(0)$. The product of binomials telescopes across the six
player steps to exactly the multinomial in \eqref{eq:Z}, so the recursion is
exact rather than approximate.

Two implementation notes make \eqref{eq:dp} fast. First, the constraint
$\sum_g \kappa_g = q_p$ is enforced with an auxiliary axis counting cards taken
so far, so the compositions $\kappa$ are never enumerated: player $p$'s step is
$O(\sum_{g \in A_p} \min(n_g,q_p))$ array operations. Second, the state space
$\prod_g (n_g+1)$ is small in practice --- median $272$, mean $577$ over real
positions --- because $G$ is small, not because $|F|$ is.

\paragraph{What the dynamic program yields.}
\begin{itemize}
\item \textbf{Partition function} $Z = |\Omega_0|$, i.e. the exact size of the
information set, computed in closed form. This makes quantities that are
normally estimated (such as the disambiguation factor of
\citealp{long2010understanding}) directly computable.
\item \textbf{Exact marginals.} A backward pass $B_p$ mirroring \eqref{eq:dp}
gives $\Ex[k_{gp}]$ by a forward-backward sweep that carries one accumulator
per group, so all $6G$ expectations come out of a single pass. Since cards
within a group are exchangeable,
$\Prob(f(c) = p) = \Ex[k_{g(c)p}] / n_{g(c)}$.
\item \textbf{Exact uniform sampling.} With $B$ available, each player's whole
take is drawn from its exact conditional, so a draw is exactly uniform on
$\Omega_0$.
\item \textbf{Exact joint probabilities.} Pinning named cards reduces $n_g$ and
$q_p$ by one each, so $\Prob(\text{a declared claim distribution is exactly
right})$ is a ratio of two evaluations of \eqref{eq:Z}. Because cards are
labelled throughout, no extra combinatorial factor is needed.
\end{itemize}

\paragraph{Validation.} On $400$ randomly generated systems the dynamic program
reproduced the brute-force count \emph{and} every per-card marginal to within
$10^{-9}$, $400/400$; the pinned-count routine matched brute force on $200/200$.
An exact uniform sampler drawn $40{,}000$ times matched the exact marginals to
$0.0000$ at four decimal places.

\begin{remark}[Completeness]
Because a marginal of zero is a proof of impossibility, the dynamic program is
a \emph{complete} propagator for the candidate-set-plus-quota subsystem: any
placement it assigns probability zero is genuinely impossible, and any
placement it assigns probability one is genuinely forced. v0.3's propagator was
explicitly documented as sound but not complete for this subsystem. The
disjunctive constraints remain outside this guarantee.
\end{remark}

% =====================================================================
\section{The disjunctive constraints}
\label{sec:clauses}

The OR-constraints are the part of the system \eqref{eq:world} that the
dynamic program does not absorb. Three routes are available, and we measured
all three rather than choosing on intuition.

\paragraph{(a) Ignore them.} On $40$ real positions, exact marginals computed
on $\Omega_0$ differed from the truth on $\Omega$ by a mean $L_1$ error of
$0.127$ per card --- \emph{worse} than v0.3's biased sampler at $0.074$. So the
clauses carry a great deal of information and cannot be dropped.

\paragraph{(b) Rejection.} Draws that are exactly uniform on $\Omega_0$ can be
filtered to be exactly uniform on $\Omega$. A typical clause covers $3.35$ live
cards, each landing with the relevant player about one time in five, so a
clause holds with probability roughly $0.53$ and five independent clauses hold
about $4\%$ of the time. Measured acceptance on real positions: mean $0.30$,
median $0.18$, and below $0.05$ on $16\%$ of positions. Rejection is therefore
exact but unaffordable.

\paragraph{(c) Fold them into the dynamic program.} Refining the groups by
OR-membership keeps the cards within a block exchangeable and makes each clause
a local condition on one player's step, which can be enforced by inclusion--
exclusion over that player's clauses. The refined state space, however, has a
bad tail: median $5{,}760$ but p90 $230{,}000$ and p99 $2.7\times10^{7}$. Exact
in the median case, hopeless in the tail.

\subsection{Sequential importance sampling with closed-form densities}

We therefore sample. Cards are assigned in a fixed order, so each world has
exactly one generating path and its proposal density $q(w)$ is the product of
the per-card choice probabilities. The proposal is quota-proportional --- the
natural approximation to uniform --- multiplied by a boost for choices that
would satisfy an as-yet-unsatisfied clause, and forced when a clause has a
single live card left. Self-normalised importance weights
\begin{equation}
\tilde{w}_i \;\propto\; \frac{L(w_i)}{q(w_i)},
\qquad
\widehat{\Ex}[\varphi] = \frac{\sum_i \tilde{w}_i \varphi(w_i)}{\sum_i \tilde{w}_i}
\label{eq:snis}
\end{equation}
are then consistent for any target $\propto L$; with $L \equiv 1$ the target is
the uniform posterior of Proposition~\ref{prop:uniform}.

Two properties are required and both are verified against exhaustive
enumeration rather than argued:

\begin{itemize}
\item \textbf{Support completeness.} Every world in $\Omega$ must have
$q(w) > 0$, since a weight cannot repair mass the sampler never reaches. We
check this \emph{analytically} for every world of every small system --- random
hit-testing cannot distinguish ``unreachable'' from ``rare'' --- over $55$
systems and $1{,}797$ worlds, all reachable.
\item \textbf{Density correctness.} The density used to build the weights must
be the density the sampler realises. Analytic and empirical path densities
agree to a worst gap of $0.0097$ over $20$ systems.
\end{itemize}

Weighted marginals then matched exhaustive enumeration to a worst error of
$0.029$ at $4{,}000$ draws, while the \emph{unweighted} estimator over the same
draws was measurably biased --- a control without which the test would not be
known to discriminate. The Kish effective sample size is $\mathrm{ESS}/N = 0.81$
on real positions, so the cost of steering the proposal is small and, more
importantly, is reported rather than assumed. Drawing the batch one card at a
time across all draws in lockstep, rather than one draw at a time, makes this
$5$--$6\times$ faster for identical output; before that change the scalar loop
was 67\% of the policy's total runtime.

\subsection{A correction: option (c) is in fact tractable}
\label{sec:clauses-correction}

Our estimate above that folding the clauses into the dynamic program is
unaffordable was measured on the wrong decomposition, and a separate study
overturned it. Refining by \emph{half-suit} is coarser than necessary: what
actually needs to stay together is a connected component of clauses that share
cards, and OR-free cards can be processed one at a time. With that
decomposition and a dynamic program over the residual-quota lattice, an exact
backend refused none of 349 real positions at a work budget of $10^7$, costing
a median of $7.5$\,ms and a p90 of $18.5$\,ms.

That result is worth stating in its own right, because it inverts the usual
shape of an accuracy-versus-compute study: \emph{the frontier here is not a
curve, it is a point}. The exact backend sits below and to the left of every
sampled one, and there is no compute budget above roughly $8$\,ms at which a
sampler is the right answer. Below that budget the best available option is a
hybrid --- exact where it fits, sampled where it does not --- which still beats
every pure sampler at matched cost ($L_1$ error $0.133$ at $5.5$\,ms against
$0.258$ for v0.3's sampler at $3.2$\,ms).

We nevertheless ship the importance sampler, and the reason is the subject of
Section~\ref{sec:oppmodel}: the exact backend computes the \emph{uniform}
posterior, and the uniform posterior is not the target we want. What the
measurements in Section~\ref{sec:res-null} add is that the difference between an
exact posterior and a well-sampled one is worth nothing measurable in play,
while the difference between the uniform target and an opponent-model target is
worth $1.9$ sets per deal-pair. Exactness was the cheaper thing to buy and the
less valuable one.

The natural synthesis is available and untried: the choice model
\eqref{eq:oppmodel} is a product of per-(player, half-suit) depth factors, and a
depth is a function of one half-suit's count vector. So if the items of the
dynamic program are taken to be whole half-suits rather than clause components,
the likelihood factorises over exactly the same structure as the constraints do,
and exact inference \emph{under the opponent model} costs the same dynamic
program with reweighted item tables. We did not build it, because the evidence
above says it would buy approximately nothing in strength; we state it because
it is the right thing to build if that evidence ever changes.

% =====================================================================
\section{Opponent modelling: relaxing the false hypothesis}
\label{sec:oppmodel}

Proposition~\ref{prop:uniform} rests on a hypothesis that is plainly false:
that a player's choice of action depends on their hand only through legality.
It does not. A player who asks in a half-suit is more likely to be deep in it.

Under the simplest defensible choice model --- a player picks which half-suit
to ask in in proportion to how many cards of it they hold ---
\begin{equation}
\Prob(\text{ask in } \halfsuit \mid \text{hand } h)
\;=\; \frac{d_\halfsuit(h)}{|h|},
\label{eq:choice}
\end{equation}
and the denominator is public, hence identical across worlds. The likelihood of
the whole observed ask sequence is therefore proportional to a product of
depths, giving the importance weight
\begin{equation}
\log L(w) \;=\; \gamma \sum_{(a,\halfsuit)} n_{a\halfsuit} \,
\log\big( d_{a\halfsuit}(w) \big),
\label{eq:oppmodel}
\end{equation}
where $n_{a\halfsuit}$ counts asks by player $a$ in half-suit $\halfsuit$ and
$d_{a\halfsuit}(w)$ is $a$'s depth in $\halfsuit$ under world $w$. The single
parameter $\gamma$ controls how sharply the model is believed;
$\gamma = 0$ recovers Proposition~\ref{prop:uniform} exactly, so the entire
idea is one ablatable number. Because \eqref{eq:oppmodel} decomposes over
cards, it costs one integer increment per card inside a draw.

\begin{remark}[The obvious refinement does not help]
Depth in \eqref{eq:oppmodel} is evaluated on the initial deal, whereas the
choice model is really about what the player held \emph{at the moment they
asked}. Cards move, so the two differ. We implemented the at-ask-time version
--- which needs only one replay of the public log per decision, not per
candidate world, because a card the propagator has not pinned was never located
at any time and so is exactly the card the sampler draws --- and it is not
better. Posterior negative log-likelihood is $1.3536$ against $1.3522$ for the
initial-deal version at matched settings, and $1.3391$ against $1.3373$ at the
best settings of each. In duplicate-deal play the two are likewise
indistinguishable. The theoretically-correct quantity buys nothing over the
simpler one, so the simpler one ships.

We flag one methodological point here because it nearly became a false finding.
Our first at-ask-time implementation counted only cards that had \emph{visibly}
changed hands, silently omitting every card the propagator had \emph{deduced} to
be someone's --- the large majority. Depths came out far too small, and the
variant measured as catastrophic on both instruments (posterior NLL $2.14$
against $1.35$; $-1.54$ sets per deal-pair in play). Two independent instruments
agreeing on a large effect is normally reassuring; here they were agreeing about
a bug. What caught it was that the magnitude was implausible for a refinement,
not that either measurement looked wrong.
\end{remark}

% =====================================================================
\section{Scoring a posterior against the truth}
\label{sec:posterior-quality}

Every belief metric in v0.3 was indirect: sampler cost, agreement with an exact
endgame solver, or downstream win rate. In simulation the hidden hands are
available, so a posterior can be scored on the question it is actually trying
to answer: how much probability does it place on where the cards really are?

We score over cards that are still genuinely uncertain for the acting seat --- a
card the propagator has already pinned is scored perfectly by every
configuration and would only dilute the comparison --- using mean negative
log-likelihood of the true holder, the six-way Brier score, top-1 accuracy, and
a reliability curve in deciles. Reliability must bin the probability assigned
to a \emph{fixed} outcome; binning by the probability assigned to the true
holder would make every entry a hit by construction and report perfect
calibration for any model whatsoever.

% =====================================================================
\section{Exact solving, and a closed form}
\label{sec:exact}

The exact solver is the project's only source of \emph{absolute} truth: in a
position small enough to solve, there is a right answer, so an engine that
disagrees has a defect rather than a preference. v0.3's solver enumerated card
placements and stopped at roughly seven live cards. Widening it produced three
results that change how the previous version's absolute benchmark should be
read.

\subsection{There is no such thing as a seven-card position}

Claims strip all six cards of a half-suit from every hand, so
\begin{equation}
\text{live cards} \;=\; 6 \times (\text{unresolved half-suits}) \qquad
\text{exactly, always.}
\end{equation}
Every ``maximum live cards'' threshold in the previous codebase was therefore a
binary switch: any value from 6 to 11 selects the same policy. v0.3 solved the
one-half-suit layer and refused everything else, and the two-half-suit layer is
$6^{12}\times 6 \approx 1.3\times10^{10}$ card placements, which is why.

\subsection{The count abstraction}

\begin{proposition}[Card identity is irrelevant under perfect information]
\label{prop:abstraction}
Let $\sigma$ be any permutation of the deck that maps each half-suit onto
itself. Then $\sigma$ is an automorphism of the perfect-information game, and
consequently the value of a position depends only on how many cards of each
live half-suit each player holds.
\end{proposition}

\begin{proof}[Proof sketch]
$\sigma$ acts on states by relabelling hands. \textsc{Ask}$(t,c)$ is legal in
$s$ iff \textsc{Ask}$(t,\sigma(c))$ is legal in $\sigma(s)$: the target and all
hand counts are unchanged, ``the asker holds a card of this half-suit'' is
preserved because $\sigma$ fixes half-suits setwise, and ``the asker does not
hold $c$'' is preserved because $\sigma$ is a bijection of each hand. The ask
succeeds iff the target holds $c$ iff $\sigma(\text{target})$ holds
$\sigma(c)$, and the transfer commutes with $\sigma$. \textsc{Claim}$(h,a)$ maps
to \textsc{Claim}$(h, a\circ\pi^{-1})$ with $\pi$ the restriction of $\sigma$ to
$h$; exactness, the ``some card sits with an opponent'' case and the
wrong-distribution case are all preserved, and both states lose exactly the six
cards of $h$. \textsc{Pass} and turn normalisation depend only on which hands
are empty. Hence $V(\sigma(s)) = V(s)$, and two states have equal
per-(half-suit, player) counts iff some such $\sigma$ relates them.
\end{proof}

A half-suit's counts are a composition of 6 into 6 ordered parts, of which there
are 462, so a layer with $m$ unresolved half-suits has at most $462^m \times 6$
abstract states: $2{,}772$ for $m=1$ and $1{,}280{,}664$ for $m=2$, against
$279{,}936$ and $1.3\times10^{10}$ card placements. Both are solved once, in
full, as dense tables --- a genuine Fish tablebase, after which a query is an
array lookup. Measured on 300 one-half-suit positions: $2{,}240$ CPU-seconds for
the old solver against $0.052$ seconds for the new one. An independent
cross-check on 131 further positions found agreement on both the value and the
complete optimal-action set in $131/131$ cases, at a $4{,}205\times$ speedup.

\subsection{A closed form for the perfect-information value}

Both full tables match, exactly, the formula
\begin{equation}
V \;=\; \mathrm{sign}(\text{mover's team}) \times (2f - m),
\label{eq:closedform}
\end{equation}
where $m$ is the number of live half-suits and $f$ the number of them in which
the team on move holds at least one card. The mechanism is that footholds only
ever shrink and a perfectly-informed asker never misses, so the team on move
takes every half-suit it can reach and the opponents take the rest.

This has an uncomfortable consequence for the previous version's headline
absolute result. v0.3 reported that its belief agents choose a provably optimal
move in 100\% of information-resolved positions --- but every one of those
positions had a single live half-suit, where \eqref{eq:closedform} reduces to
``the team on move wins it''. The score was real, and it is reproduced here, but
it is a shallower certificate than it reads. The two-half-suit slice is the one
with teeth, and the corpus used in Section~\ref{sec:res-absolute} is 79\%
two-half-suit for that reason.

\subsection{Bellman-optimal is not optimal in a loopy game}

Fish's state graph is cyclic, so the Bellman operator has many fixpoints and
value-preservation is not optimality: if a side can force $+1$, every action
keeping the value at $+1$ looks optimal, including one that walks a cycle
forever and banks nothing. Measured on the full tables, in \textbf{28.6\%} of
decided two-half-suit states some value-preserving action makes no progress. At
one half-suit the figure is $0\%$, which is why v0.3 never noticed. The solver
therefore carries an attractor rank and the agent asks for a
\emph{progress-optimal} action, and Section~\ref{sec:res-absolute} reports
agreement against that stricter criterion.

Two further corrections to v0.3. First, although the state graph is cyclic,
\emph{optimal play always terminates} under perfect information: cycles exist,
but no optimal line needs one. Second, the fixpoint reached by value iteration
was checked to be unique in real Fish layers --- eight different initialisations
(zeros, $\pm m$, $\pm 5$, three random) and three sweep orders at $m=1$ and
$m=2$ all land on the attractor value --- so v0.3's zero-initialisation and
dictionary ordering were never load-bearing, even though in general they could
have been.

% =====================================================================
\section{The ask objective}
\label{sec:objective}

\subsection{A wider basis}

v0.3 scored an ask as $\Prob(\text{success}) + 0.06 \cdot \text{depth}$, later
adding a turn-risk and a scarcity term. Those terms are in incommensurable
units --- a probability, a card count, a normalised hand-size difference, a
share --- so the weights combining them have no meaning beyond ``what won a
sweep''. That is consistent with what v0.3 measured: a clean inverted-U in each
weight, which is the signature of a \emph{tie-breaker} rather than an objective.

v0.4 keeps the additive form, because it is ablatable and because the previous
version's evidence says the success probability really is dominant, but widens
the basis to eleven terms and computes all of them from the posterior in a
single vectorised pass. Writing $\pi$ for $\Prob(\text{success})$ and $H$ for
the asked card's half-suit, the score is
\begin{equation}
\text{score}(a) \;=\; \pi \;+\; \sum_{j} w_j\, f_j(a),
\label{eq:objective}
\end{equation}
with the terms in Table~\ref{tab:terms}. Two deserve comment because they are,
as far as the surveyed literature shows, unstudied.

\paragraph{Exposure.} v0.3's turn-risk term penalised handing the turn to an
opponent with many cards. Hand size is a proxy; the quantity that actually
matters is how much that opponent can \emph{immediately take from us}, which we
can compute: cards of ours whose location is already \emph{public} and that sit
in a half-suit the opponent can legally ask in. Both parts are public knowledge,
so using them is inference, not leakage.

\paragraph{Signalling.} An ask publicly certifies that the asker holds a card of
that half-suit. In a game with no legal communication channel, that is the only
channel there is --- and it is simultaneously read by the opposition. v0.3
modelled only the cost (its \texttt{reveal} term, worth a marginal $+0.29$). The
benefit is that a teammate learns where a card of a half-suit our team is
winning must be, which is exactly the localisation that lets a set be declared.
So the sign should depend on who gains more, and the natural proxy is our team's
share of the half-suit. Cost and benefit are carried as independent weights so
that the data, not the story, decides.

\begin{table}[t]
\centering
\small
\begin{tabular}{llp{7.4cm}}
\toprule
term & origin & feature value \\
\midrule
\texttt{suit}    & v0.3 & own depth in $H$ (raw card count) \\
\texttt{turn}    & v0.3 & $(1-\pi) \cdot$ negated normalised hand-size excess of the target \\
\texttt{scarce}  & v0.3 & $2(\bar{s}_H - \tfrac12)$, with $\bar{s}_H$ our team's expected share of $H$ \\
\texttt{reveal}  & v0.3 & $-1$ if $H$ has not previously been shown by us \\
\texttt{deplete} & v0.3 & $\pi \cdot (1 - \text{target hand}/q)$ \\
\texttt{expose}  & new  & $-(1-\pi) \cdot$ our publicly located cards the target could take \\
\texttt{claim}   & new  & $\pi \cdot \big(\prod_{c \in H}\Prob(\text{$c$ with our team})\big)^{1/6}$ \\
\texttt{info}    & new  & $4\pi(1-\pi) \cdot$ holder entropy of $H$ \\
\texttt{certain} & new  & $1$ if $\pi = 1$ \\
\texttt{concent} & new  & largest single teammate share of our team's holding in $H$ \\
\texttt{signal}  & new  & $2(\bar{s}_H - \tfrac12)$ if $H$ has not previously been shown by us \\
\bottomrule
\end{tabular}
\caption{The ask objective basis. Every term is zero by default, so setting a
weight to zero ablates exactly that idea and nothing else --- a property
enforced by test, for the reason given in Section~\ref{sec:eval}.}
\label{tab:terms}
\end{table}

\subsection{An alternative in the units that are actually paid out}

The additive form has a principled competitor. What the game pays is half-suits,
so define, for each unresolved half-suit,
\begin{equation}
V(H) \;=\; \Prob(\text{our team ends up scoring } H) - \Prob(\text{they do}),
\label{eq:hsvalue}
\end{equation}
which is measured in sets, is bounded in $[-1,1]$, and whose sum over half-suits
is the final differential. An ask is then worth the expected change it produces,
\begin{equation}
\text{score}(a) \;=\; \pi\,\big[V(H \mid \text{success}) - V(H)\big]
\;+\; (1-\pi)\,\big[V(H \mid \text{failure}) - V(H)\big],
\label{eq:valobjective}
\end{equation}
with every quantity in the same unit. The substantive difference is that this is
a \emph{marginal} value rather than a level: ``our team is winning this
half-suit'' is a level, and what should drive a decision is how much one more
card moves the outcome.

$V$ is obtained by supervised learning rather than by guessing a functional form:
self-play games supply, for every unresolved half-suit at every decision, a
feature vector computable from the acting seat's posterior alone, labelled with
how that half-suit eventually resolved. That is roughly nine labels per
decision, which is abundant and cheap, in contrast to the rollout regression
that v0.3's variance analysis made expensive. The counterfactual posteriors in
\eqref{eq:valobjective} are local updates on the half-suit's marginal block: a
success pins the asked card to us, a failure removes the target from that card's
candidate set.

% =====================================================================
\section{Claiming}
\label{sec:claiming}

v0.3 swept the claim-confidence threshold directly and found two things. First,
claiming eagerly is a real, measurable mistake: dropping the bar from $0.97$ to
$0.70$ costs $0.45$ sets per deal-pair. Second, and more usefully, thresholds
from $0.85$ to near-certainty play \emph{identically}: over 150 duplicate deals
the $0.97$ and $0.999$ policies never once diverged.

The mechanism is worth stating, because it justifies the design rather than
leaving $0.97$ as a magic number. Suppose our team genuinely holds all six cards
of a half-suit but we cannot place them within the team. What does waiting cost?
An opponent cannot take the set by claiming it --- a claim by a player whose team
does not hold all six \emph{gives} the set to us. So the only cost of waiting is
running out of opportunity, and the benefit is that continued play localises
teammate holdings. Waiting is close to free, which is exactly what the sweep
measured and what v0.3's expected-value model missed by treating the resolution
of a split as a coin flip.

The corollary is that the leverage in claiming is not the threshold. It is
(i)~getting the declared distribution right, since nulls cost roughly $0.6$ sets
per deal-pair and a null is a bookkeeping failure rather than an unlucky gamble;
and (ii)~playing the forced claims well, where no threshold is involved at all.
v0.4 therefore keeps a high threshold and changes what happens underneath it:

\begin{itemize}
\item The declared distribution is the posterior MAP, found by shortlisting on
the per-card marginals and scoring the shortlist with the \emph{joint}
posterior. The product of marginals is not the joint --- cards compete for the
same quota slots, so per-card modes can be jointly impossible.
\item Where the clause set is empty the joint probability is computed exactly by
the dynamic program of Section~\ref{sec:counting}; where it is not, deduction
still settles the extremes exactly, because a zero under the OR-free system
implies a zero under the smaller OR-constrained one, and likewise for one.
\item Forced claims maximise expected sets, $\Prob(\text{exact}) -
\Prob(\text{an opponent holds one})$, rather than confidence.
\end{itemize}

Finally, a hard architectural constraint taken from the literature: a claim must
never be selected by determinized search. Inside any determinization the
searcher knows the true layout, so every claim evaluates as a certain gain and
the declared assignment differs between sampled worlds --- textbook strategy
fusion \citep{frank1998search}. Claims in v0.4 are always decided from the
posterior.

% =====================================================================
\section{Search over the belief}
\label{sec:lookahead}

Every search this project has tried --- PIMC, Information-Set MCTS, and two
value-network variants --- evaluates a candidate move by sampling hidden layouts
and solving each one, and every one of them lost to its own prior. v0.3 did not
leave that as a mystery: it measured the standard deviation of a position's
value across possible layouts at about $2.4\times$ the gap between the best and
the worst candidate move, so a searcher scoring different candidates against
different sampled worlds is ranking sampling noise. Common random numbers
removed the deficit and never produced a surplus.

That diagnosis rules out a family of designs, not search itself. What it rules
out is sampling. So the design it leaves open is a lookahead whose state is the
\emph{posterior} rather than a sample from it: the marginal matrix
$M[c,p] = \Prob(\text{player } p \text{ holds card } c)$, with transitions that
are local edits of that matrix. Given a belief such a search is a deterministic
function --- run it twice on one position and it returns the same vector --- so
the variance that killed the previous attempts is structurally absent rather
than merely reduced.

\subsection{What it computes}

A Literature turn is a run of asks that ends the instant one misses. Greedy
scoring maximises the immediate success probability; what a possession is
actually worth is how many cards it banks before the turn flips. Define
\begin{equation}
W(B, d) \;=\; \max_{a \,\in\, \mathcal{A}(B)}\;
p_a \,\bigl[\, 1 + W(B \mid a \text{ succeeded},\, d-1) \,\bigr],
\qquad W(B, 0) = 0,
\label{eq:chain}
\end{equation}
where $\mathcal{A}(B)$ is the legal-ask set implied by the belief and
$p_a = M[\text{card}(a), \text{target}(a)]$. The failure branch contributes
nothing to $W$ because a miss ends the possession; what a miss costs
positionally is already priced by the incumbent objective's turn-risk and
exposure terms, and duplicating it here would double-count.

$W$ is in units of cards expected to be banked in the rest of this possession.
It enters the policy as an additive bonus, not as a replacement:
\begin{equation}
\text{score}(a) \;=\; \underbrace{\pi_a + \textstyle\sum_j w_j f_j(a)}
_{\text{the incumbent objective, unchanged}}
\;+\; w_{\mathrm{look}} \cdot p_a \, W(B \mid a \text{ succeeded},\, d-1).
\label{eq:lookscore}
\end{equation}
The reason it is additive is Section~\ref{sec:res-null}: replacing
$\Prob(\text{success})$ with a learned half-suit value --- correct units, a
calibrated model --- costs $7.4$ sets per deal-pair. Whatever else this game
rewards, the probability that the ask lands is the objective and not a term in
one. At $d \le 1$ the bonus is identically zero, so the policy reproduces the
incumbent decision for decision; that is enforced by test, for the reason
Section~\ref{sec:eval} gives.

\subsection{Why the obvious version is provably a no-op}

The natural transition is to collapse the asked card's row to a point mass on
us. That is exact, and on its own it is \emph{inert}, which is worth stating
because it determines whether the design can do anything at all.

\begin{proposition}[Greedy optimality under an inert transition]
\label{prop:greedy}
Suppose the transition edits only the asked card's row, so that $p_{a'}$ is
unchanged for every $a' \ne a$, and suppose every candidate remains available
after any other is taken. Then \eqref{eq:chain} is maximised by taking
candidates in descending order of $p$, and its value equals that of the greedy
policy at every depth.
\end{proposition}

\begin{proof}[Proof sketch]
Under the hypotheses, a depth-$d$ chain is a choice of an ordered $d$-subset of
a \emph{fixed} multiset of probabilities, with value
$p_{1}(1 + p_{2}(1 + \dots))= \sum_{k=1}^{d} \prod_{i \le k} p_{i}$. Exchanging
an adjacent pair $(p_i, p_{i+1})$ with $p_i < p_{i+1}$ changes the value by
$\bigl(\prod_{j<i} p_j\bigr)(p_{i+1} - p_i) \ge 0$, so descending order is
optimal; and the greedy policy selects exactly that order.
\end{proof}

Proposition~\ref{prop:greedy} is the null this section is tested against, and it
is why a lookahead built only on the point-mass update would be worth measuring
only as a bug-check. Its hypotheses fail in three places, of which the third is
the substantive one:

\begin{enumerate}
\item a card taken cannot be asked for again, and one card is askable from
several targets, so candidates do disappear in correlated groups;
\item a target can run out of cards, which removes every candidate aimed at
them;
\item \textbf{the quota system couples the candidates.} Learning that the target
held this card makes it less likely they hold the others --- they are now known
to hold one card fewer, and that probability mass has to go somewhere.
\end{enumerate}

The third is handled by scaling the target's column down to the count they now
have and renormalising each affected row across the remaining players: one sweep
of the same proportional fitting the exact posterior does globally. A row whose
whole mass sits on the shrunk column renormalises straight back to where it
started, which is correct --- a publicly located card cannot be made less
certain by a counting argument. The sweep costs one column scale and one row
normalisation per node, and it is separately ablatable, so
Proposition~\ref{prop:greedy} is not merely argued but measured: the
no-coupling cell in Table~\ref{tab:lookahead} is the experiment.

\subsection{Cost}

The branching factor of a real mid-game position is large: measured over 615
decisions from six self-play games at the champion configuration, the number of
legal asks has mean $45.0$, median $45$, p90 $81$ and maximum $102$, and exceeds
$8$ in $94.6\%$ of decisions. (An earlier draft of this section said the
branching factor ``averages 8''. That figure is the mean legal-action count of
the \emph{exact-solver corpus} in
\texttt{results/exact\_positions\_v2.json}, which consists of endgame positions
with one or two live half-suits --- an endgame statistic read as a general one.)

Continuations are beamed to the best $4$ candidates by immediate $p$. The root
ply is \emph{not} beamed, because the policy needs a bonus for every candidate it
is ranking, so a naive depth-3 tree costs $n\,(1 + b + b^2) = 21n$ expansions ---
about $945$ at the median $n = 45$ --- rather than the $b + b^2 + b^3 = 84$ a
uniformly-beamed tree would give.

Three quarters of those expansions were doing no work. At the last ply the
continuation is $W(\cdot, 0) = 0$ by definition, so $q = p\,(1 + 0) = p$ and the
maximum over the beam is simply the largest probability --- a number the sort key
has already computed. The implementation nonetheless performed up to $b$ full
transition/undo cycles there, each copying the marginal matrix and running a
column scale and a $54$-row renormalisation, to rediscover it. Short-circuiting
that ply, and fusing the row renormalisation to avoid a boolean-index round trip,
leaves the returned vector bit-identical (maximum absolute difference $0.0$ over
$24$ real positions and every candidate) and costs $n\,(1 + b) = 225$ expansions
instead of $945$. Measured median cost falls from about $14$\,ms per decision to
$6.7$\,ms (mean $8.3$, p90 $18.1$, over $40$ positions, minimum of three
repetitions each) on the contended host of Section~\ref{sec:counting}.

Against the configuration the lookahead runs actually used --- $160$ draws, the
\texttt{fish4/agent4.py} default, since no lookahead cell in
\texttt{results/v04\_duels.jsonl} overrides \texttt{n\_draws} --- the relevant
posterior cost is the $31.9$\,ms row of Table~\ref{tab:posterior}, not the
$\sim\!100$\,ms rows, which are all $512$-draw. So the bonus now adds roughly a
fifth to a decision. An earlier draft of this section claimed $15\%$ by comparing
the pre-optimisation cost against a $512$-draw budget these runs never used,
which was wrong twice in opposite directions; the duelling results in
Table~\ref{tab:lookahead} were all produced at the slower $945$-expansion cost,
since the optimisation is exact and postdates them.

% =====================================================================
\section{Evaluation methodology}
\label{sec:eval}

\subsection{Duplicate deals}

Raw win rates in a card game are dominated by the deal. Every result here uses
\textbf{duplicate (paired) deals}: each deal is played twice with the teams
swapped, on identical cards, an identical rotated starting seat, and identical
agent randomness. Only the policy assignment differs, so per-pair set
differentials are i.i.d.\ across deals and confidence intervals mean what they
say.

\subsection{How much can an experiment of size \texorpdfstring{$n$}{n} resolve?}

The per-pair standard deviation of the set differential between two policies of
the strength we compare (v0.3's champion against its own baseline) is
\textbf{3.869 sets}, measured over 300 pairs. The distribution is wide: the
theoretical bound is $\pm 18$ and observed pairs stayed within $\pm 11$, but a
single pair says almost nothing. The consequences are stark:

\begin{center}
\small
\begin{tabular}{lrrr}
\toprule
effect (sets/pair) & pairs @ 80\% power & pairs @ 90\% & hours @ 80\% \\
\midrule
0.10 & 11\,751 & 15\,731 & 3.9 \\
0.20 & 2\,938 & 3\,933 & 0.98 \\
0.30 & 1\,306 & 1\,748 & 0.44 \\
0.50 & 471 & 630 & 0.16 \\
1.00 & 118 & 158 & 0.04 \\
\bottomrule
\end{tabular}
\end{center}

Read the other way: the 600-pair cells of the v0.3 champion search could resolve
effects down to about $0.44$ sets per pair, so every ``no significant
difference'' it reported below that size was \emph{uninformative rather than
negative}. We therefore report, alongside every interval, the smallest effect
the run could have detected.

\subsection{Anytime-valid stopping}

Screening many candidate terms invites peeking. A fixed-$n$ 95\% $t$-interval,
monitored after every deal, rejects a \emph{true null} 68.8\% of the time on an
empirical null and 92.1\% on a skewed one, with a median false stop after 11 and
3 deals respectively. We therefore use anytime-valid procedures: an
empirical-Bernstein confidence sequence and a betting test supermartingale
\citep{waudbysmith2024estimating}. Measured type-I error under continuous
monitoring over $4000$ null experiments each monitored at every one of $4000$
deals: 2.1--2.4\% for the betting test and 0.3--0.5\% for the confidence
sequence, against a nominal 5\%.

\subsection{Two corrections to the previous harness}

\paragraph{Timed-out games must be scored, not dropped.} Nothing in Literature
forces progress, so a harness needs an action cap. v0.3 dropped any deal-pair
containing a capped game. Dropping conditions on the outcome, which is precisely
what a paired design exists to avoid, and an audit conducted for this paper
found it biases rather than merely blunts: replaying the same 60 deals at
reduced caps, the dropped pairs had a true mean differential of $-7.14$ against
$-10.13$ for the kept ones, and at one cap the reported 95\% interval
\emph{excluded} the true value. v0.4 scores a capped game on the half-suits
actually resolved, with the unresolved ones counting for nobody --- the same
semantics the exact solver gives an unbroken cycle.

\paragraph{An ablation must reduce to its baseline.} v0.3's most instructive
methodological failure was a confounded ablation: the ablated agent silently
normalised a second term differently, so every cell compared two changes at once
and produced tight, confident, and completely wrong intervals. The confound was
caught only by noticing an implausible monotonic pattern in the magnitudes. We
therefore enforce by test that a zero weight reproduces the baseline
\emph{decision for decision} over real positions, and separately that a non-zero
weight demonstrably changes some decision. Every term in
Table~\ref{tab:terms} passes both.

% =====================================================================
\section{Results}
\label{sec:results}

\subsection{Exact inference alone changes nothing}
\label{sec:res-exact}

The first experiment is the one the research programme predicted would pay.
Holding the ask objective at v0.3's champion weights and changing \emph{only}
the inference --- exact counting where the clause set is empty, unbiased
importance sampling where it is not, in place of a sampler documented as
non-uniform --- gives a set differential of $-0.22$ per deal-pair,
95\% CI $[-0.96, +0.52]$, over 100 duplicate deals. Statistically
indistinguishable, and the interval is wide enough that only effects above
about $1.1$ sets per pair could have been detected.

This is a genuine negative result about a change that is provably correct. The
explanation is in the next subsection.

\subsection{Posterior quality against ground truth}
\label{sec:res-posterior}

Scoring posteriors against the true hidden hands over 517 decisions and 13{,}640
uncertain-card predictions (Table~\ref{tab:posterior}) shows what happened.

\begin{table}[t]
\centering
\small
\begin{tabular}{lrrrr}
\toprule
inference & NLL $\downarrow$ & Brier $\downarrow$ & top-1 $\uparrow$ & ms/decision \\
\midrule
exact DP, clauses ignored          & 1.4776 & 0.7568 & 0.287 & 50.4 \\
v0.3 sampler, 32 draws             & 1.4399 & 0.7065 & 0.386 & 5.9 \\
v0.4 importance sampler, 160       & 1.3731 & 0.7006 & 0.396 & 29.5 \\
v0.4 importance sampler, 512       & 1.3618 & 0.6955 & 0.403 & 92.3 \\
v0.3 sampler, 96 draws             & 1.3552 & 0.6891 & 0.406 & 16.6 \\
\textbf{v0.3 sampler, 512 draws}   & \textbf{1.3408} & 0.6833 & 0.416 & 87.8 \\
v0.4 + opponent model, 160, $\gamma{=}0.45$ & 1.3476 & 0.6856 & 0.401 & 31.9 \\
v0.4 + opponent model, 512, $\gamma{=}0.30$ & 1.3332 & 0.6791 & 0.415 & 99.8 \\
\textbf{v0.4 + opponent model, 512, $\gamma{=}0.45$} & \textbf{1.3251} & \textbf{0.6747} & 0.413 & 98.6 \\
v0.4 + opponent model, 512, $\gamma{=}0.60$ & 1.3248 & 0.6744 & 0.411 & 100.8 \\
\bottomrule
\end{tabular}
\caption{How much probability each inference configuration places on where the
cards actually are. Lower NLL and Brier are better. The row that matters is the
comparison between the two bold entries: the \emph{exactly correct} uniform
posterior at 512 draws (1.3618) is worse than v0.3's \emph{biased} sampler at
the same budget (1.3408).}
\label{tab:posterior}
\end{table}

Two things stand out.

First, the exactly-uniform posterior is \emph{worse} at predicting reality than
the heuristic sampler it replaced: $1.3618$ against $1.3408$ at a matched draw
budget. That is not a contradiction. Proposition~\ref{prop:uniform} is a correct
theorem about a false hypothesis. v0.3's sampler seeded the OR-clauses by
picking a live card of the asked half-suit for the asker, which systematically
over-weights worlds in which askers are deep in the suits they asked in --- an
accidental opponent model. Removing the bias removed the model with it.

Second, making the model explicit and tunable recovers the loss and passes both:
$1.3251$ at $\gamma = 0.45$. And the improvement is not a large-budget artefact.
At 160 draws the opponent model takes the sampler from $1.3731$ to $1.3476$,
which is better than v0.3's deployed 32-draw configuration ($1.4399$) by a wide
margin.

Third, and separately: ignoring the OR-clauses is by far the worst option
tested, confirming the design decision of Section~\ref{sec:clauses} from the
other direction.

\subsection{The opponent model in play}
\label{sec:res-oppmodel}

Posterior accuracy is not strength. Table~\ref{tab:gamma} measures the model by
duplicate-deal play against the otherwise-identical $\gamma = 0$ policy.

\begin{table}[t]
\centering
\small
\begin{tabular}{lrrl}
\toprule
$\gamma$ & pair score & set diff / pair & 95\% CI \\
\midrule
0.10 & 0.613 & $+0.993$ & $[+0.470, +1.517]$ \\
0.30 & 0.667 & $+1.847$ & $[+1.230, +2.463]$ \\
0.45 & 0.700 & $+1.860$ & $[+1.336, +2.384]$ \\
0.60 & 0.657 & $+1.300$ & $[+0.795, +1.805]$ \\
0.90 & 0.723 & $+1.827$ & $[+1.199, +2.455]$ \\
\bottomrule
\end{tabular}
\caption{Opponent model against the identical policy with the model disabled,
150 duplicate deal-pairs per cell. Every cell excludes zero. The dose-response
rises to a plateau by $\gamma = 0.30$ and stays there over a threefold range,
which is what one wants from a model with a single free parameter --- and unlike
the ask-objective terms of v0.3, which showed a sharp inverted-U, there is no
weight at which this becomes destructive within the range tested.}
\label{tab:gamma}
\end{table}

\paragraph{The cleanest attribution.} Tables~\ref{tab:gamma} measures the model
against copies of itself. The sharper experiment holds the \emph{opponent} fixed
at the v0.3 champion and varies only whether the model is on, at identical deal
seeds and identical agent seeds, 250 duplicate deal-pairs each:

\begin{center}
\small
\begin{tabular}{lrrl}
\toprule
v0.4 configuration & pair score & set diff / pair & 95\% CI \\
\midrule
exact inference, $\gamma = 0$ & 0.502 & $-0.008$ & $[-0.482, +0.466]$ \\
exact inference, $\gamma = 0.30$ & 0.684 & $\mathbf{+1.920}$ & $[+1.479, +2.361]$ \\
\bottomrule
\end{tabular}
\end{center}

The first row is as close to exactly zero as a 250-pair experiment can report.
Provably correct inference, replacing a sampler its own authors documented as
biased, moved the result by eight thousandths of a set per deal-pair. The second
row differs from the first in one parameter.

\subsection{The strength ladder}
\label{sec:res-ladder}

\begin{table}[t]
\centering
\small
\begin{tabular}{lrrll}
\toprule
v0.4 champion versus & pairs & pair score & set diff / pair & 95\% CI \\
\midrule
\textbf{v0.3 champion (\texttt{tuned})} & \textbf{500} & \textbf{0.705} & $\mathbf{+1.854}$ & $\mathbf{[+1.542, +2.166]}$ \\
\quad --- replication & 300 & 0.758 & $+2.440$ & $[+2.052, +2.828]$ \\
\quad --- replication, $\gamma = 0.30$ & 250 & 0.684 & $+1.920$ & $[+1.479, +2.361]$ \\
v0.3 \texttt{probabilistic}    & 200 & 0.792 & $+2.785$ & $[+2.274, +3.296]$ \\
\bottomrule
\end{tabular}
\caption{The v0.4 champion is
\texttt{fishbot4($\gamma$=0.35)}: exact/unbiased inference, the opponent choice
model, and v0.3's own ask weights. No new objective term is enabled.}
\label{tab:ladder}
\end{table}

Against the two weakest rungs the champion wins every pair: $120/120$ against
the public-information heuristic ($+14.09$ sets per pair) and $80/80$ against
random ($+16.30$). Those are reported for completeness rather than as evidence;
a shutout means the rating gap is a bound, not a measurement.

The headline comparison was run three times at different seeds, and the three
estimates span $+1.85$ to $+2.44$ --- wider than their individual intervals
suggest, which is a useful reminder that a 95\% interval covers the mean of the
deals actually played, and that three such intervals will occasionally fail to
overlap. We take the largest experiment as the headline and report the others
rather than the most flattering one.

\paragraph{Do v0.3's objective terms survive better inference?} A reasonable
worry about the whole result is that v0.3's hand-tuned tie-breakers were
compensating for a poor posterior and would become redundant once it improved.
They do not become redundant, but they do become \emph{worth less}. Removing
both --- setting $w_{\mathrm{turn}}$ and $w_{\mathrm{scarce}}$ to zero and
changing nothing else --- costs, over 250 duplicate deals each:

\begin{center}
\small
\begin{tabular}{lrl}
\toprule
inference & cost of removing both terms & 95\% CI \\
\midrule
without the opponent model & $1.660$ sets/pair & $[+1.240, +2.080]$ \\
with the opponent model    & $1.084$ sets/pair & $[+0.611, +1.557]$ \\
\bottomrule
\end{tabular}
\end{center}

v0.3 measured the same pair at $+1.41$ against its own inference, which sits
between the two. So the terms and the opponent model are partially substitutable:
about a third of what the tie-breakers were contributing is absorbed once the
posterior knows more about opponents' holdings, and the remaining two-thirds is
positional judgement that no amount of card-reading supplies. That is the
expected shape if the tie-breakers were doing two jobs --- compensating for a
blurry posterior, and encoding something genuinely strategic --- and it is worth
noting that the gains do \emph{not} simply add: $1.85$ (the headline) is less
than $1.92$ (the model alone against the v0.3 champion) plus anything.

\subsection{Absolute strength: agreement with provably optimal play}
\label{sec:res-absolute}

Every number above is relative. Table~\ref{tab:exact} is not. It reports, over
988 exactly solved positions harvested from 260 real games by four different
generator policies, how often each agent chooses a \emph{progress-optimal}
action --- the stricter criterion of Section~\ref{sec:exact}, which excludes
value-preserving moves that make no progress.

\begin{table}[t]
\centering
\small
\begin{tabular}{lrrrr}
\toprule
agent & resolved & uncertain & $m{=}1$ & $m{=}2$ \\
\midrule
random             & 60.8\% & 32.7\% & 28.3\% & 43.9\% \\
heuristic          & 71.6\% & 46.8\% & 34.4\% & 59.0\% \\
memory             & \textbf{100.0\%} & 61.7\% & 81.1\% & 70.1\% \\
probabilistic      & \textbf{100.0\%} & 66.9\% & 86.8\% & 73.3\% \\
v0.3 champion      & \textbf{100.0\%} & 67.6\% & 86.8\% & 74.0\% \\
v0.4, no opponent model & \textbf{100.0\%} & 67.5\% & 85.4\% & 74.2\% \\
\textbf{v0.4 champion}  & \textbf{100.0\%} & \textbf{72.5\%} & 85.8\% & \textbf{78.7\%} \\
v0.4 champion, tablebase disabled & \textbf{100.0\%} & 72.5\% & 85.8\% & 78.7\% \\
\bottomrule
\end{tabular}
\caption{Agreement with progress-optimal play over 988 solved positions, 278 of
them information-resolved and 776 with two live half-suits. ``resolved'' and
``uncertain'' mean what they did in v0.3; $m$ is the number of live half-suits.}
\label{tab:exact}
\end{table}

Four things are worth drawing out.

\textbf{Every belief-tracking agent is progress-optimal in every
information-resolved position}, on a corpus four times the size of v0.3's and
79\% of it two-half-suit, and against the stricter criterion. This reproduces
v0.3's strongest correctness claim on harder ground --- while
Section~\ref{sec:exact} explains why the one-half-suit version of that claim was
shallower than it read.

\textbf{The absolute metric confirms the relative one, independently.} Exact
inference without the opponent model scores $67.5\%$ under uncertainty against
the v0.3 champion's $67.6\%$: no change, exactly as the duplicate-deal
experiment found and exactly as the ground-truth posterior scores predicted.
Adding the opponent model moves it to $72.5\%$. Three different instruments ---
paired play, posterior likelihood against the true hands, and agreement with an
exact solver --- agree on which component does the work.

\textbf{The tablebase contributes nothing here}, and we report that rather than
claiming the speedup as strength. Disabling it changes not one of the 988
decisions. In an information-resolved position the marginals are one-hot and the
ordinary policy already plays it correctly; in an uncertain position the
tablebase refuses by construction. Its value is a $4{,}205\times$ reduction in
the cost of the endgame path and insurance against a class of error v0.3
demonstrably suffered, not extra strength.

\textbf{Where the remaining gap is.} At $m{=}2$ under uncertainty the champion
agrees with optimal play $78.7\%$ of the time. The missing fifth is not in the
endgame in the sense v0.3 meant --- these positions \emph{are} the endgame --- it
is in choosing correctly when the hidden information genuinely has not resolved.

\subsection{Everything else we tried}
\label{sec:res-null}

Table~\ref{tab:nulls} is the rest of the study. It is mostly nulls, and it is
reported in full because a table of only the things that worked would
misrepresent how much was tried. Each cell is a duplicate-deal duel against the
champion with exactly one thing changed, and the last column is the smallest
effect that cell could have resolved at 80\% power given the measured per-pair
standard deviation of 3.869 sets.

\begin{table}[t]
\centering
\small
\begin{tabular}{llrrl}
\toprule
what changed & setting & pairs & set diff & 95\% CI \\
\midrule
\multicolumn{5}{l}{\emph{inference precision}}\\
draws per decision & 96 (from 160) & 160 & $-0.300$ & $[-0.852, +0.252]$ \\
draws per decision & 320 (from 160) & 200 & $-0.035$ & $[-0.578, +0.508]$ \\
count mode & $\sqrt{\cdot}$, $\gamma{=}0.70$, 160 & 200 & $+0.480$ & $[-0.042, +1.002]$ \\
count mode & $\sqrt{\cdot}$, $\gamma{=}0.70$, 320 & 200 & $-0.165$ & $[-0.701, +0.371]$ \\
no-declaration signal & $\lambda = 0.9$ & 200 & $+0.205$ & $[-0.301, +0.711]$ \\
no-declaration signal & $\lambda = 1.8$ & 200 & $+0.190$ & $[-0.313, +0.693]$ \\
\midrule
\multicolumn{5}{l}{\emph{new ask-objective terms}}\\
exposure & $0.40$ & 160 & $-0.113$ & $[-0.660, +0.435]$ \\
claim progress & $0.30$ & 160 & $-0.206$ & $[-0.807, +0.395]$ \\
signalling & $+0.25$ & 160 & $-0.300$ & $[-0.890, +0.290]$ \\
signalling & $-0.25$ (conceal) & 160 & $+0.056$ & $[-0.530, +0.642]$ \\
information gain & $0.10$ & 160 & $-0.394$ & $[-0.932, +0.144]$ \\
concentration & $0.15$ & 160 & $-0.037$ & $[-0.653, +0.578]$ \\
turn-risk retune & $0.35$ (from $0.60$) & 160 & $-0.075$ & $[-0.670, +0.520]$ \\
scarcity retune & $0.35$ (from $0.20$) & 160 & $+0.125$ & $[-0.381, +0.631]$ \\
\midrule
\multicolumn{5}{l}{\emph{claiming and the endgame}}\\
claim threshold & $0.99$ (from $0.97$) & 200 & $+0.010$ & $[-0.010, +0.030]$ \\
informed pass choice & on & 160 & $+0.013$ & $[-0.012, +0.037]$ \\
tablebase & one half-suit only & 200 & $+0.005$ & $[-0.005, +0.015]$ \\
tablebase & disabled entirely & 200 & $+0.005$ & $[-0.005, +0.015]$ \\
\midrule
\multicolumn{5}{l}{\emph{the learned half-suit value objective}}\\
blended & $w_{\mathrm{value}} = 0.5$ & 200 & $+0.410$ & $[-0.086, +0.906]$ \\
blended & $w_{\mathrm{value}} = 1.5$ & 200 & $+0.040$ & $[-0.502, +0.582]$ \\
\textbf{replacing the objective} & pure $\Delta V$ & 200 & $\mathbf{-7.355}$ & $[-7.875, -6.835]$ \\
\bottomrule
\end{tabular}
\caption{Everything else. Only the last row resolves, and it resolves against
us. Every other interval contains zero at a minimum detectable effect of
roughly $0.8$ sets per pair, so these are ``not shown to matter at this budget'',
not ``shown not to matter''.}
\label{tab:nulls}
\end{table}

Four of these deserve comment.

\textbf{Precision has saturated, and that is new.} v0.3 found that raising the
sampled-world count from 32 to 96 was worth $+0.54$ sets per pair and concluded
that belief precision had not saturated. Under v0.4's inference it plainly has:
halving or doubling the draw budget around 160 changes nothing measurable. The
natural reading is that v0.3 was buying accuracy, not precision --- more draws
partially averaged out a \emph{bias} --- and once the bias is gone, the
remaining sampling noise is no longer the binding constraint.

\textbf{One cell technically resolves, and should still be ignored.} Lowering
the claim bar from $0.97$ to $0.90$ gives $+0.035$ $[+0.007, +0.063]$: an
interval excluding zero, and an effect of one fortieth of a set per deal-pair.
After forty-odd cells at a nominal 95\% level, an interval that narrowly
excludes zero at a magnitude this small is not evidence of anything, and we
report it rather than promoting it --- particularly having already watched a
larger apparent effect fail to replicate twice
(Section~\ref{sec:res-replication}). The neighbouring cell in the other
direction ($0.99$) gives $+0.010$ $[-0.010, +0.030]$, which is the shape of a
policy that barely changes its decisions at all.

\textbf{The claim threshold and the tablebase are exact no-ops.} The intervals
$[-0.010, +0.030]$ and $[-0.005, +0.015]$ are not small effects; they are almost
no divergence at all. Raising the claim bar from $0.97$ to $0.99$ changes
essentially no decisions, reproducing v0.3's bimodality finding on a much better
posterior. Disabling the exact endgame solver entirely changes essentially
nothing, which agrees exactly with the absolute benchmark
(Section~\ref{sec:res-absolute}), where it changed not one of 988 decisions.

\textbf{Replacing the objective with the learned value is catastrophic.} The
half-suit value model is good --- held-out log-loss $0.727$ against $0.821$ for
the class prior and $0.743$ for the best single feature, mean calibration bias
$-0.006$, reliable across every decile. Scoring asks by the expected change in
it, in units of sets, loses $7.4$ sets per deal-pair. Blending it in at low
weight is harmless and does nothing. This is the sharpest confirmation of v0.3's
central claim about this game that we have: the probability that an ask lands is
not one term among several, it is the objective, and a well-calibrated model of
what a half-suit is worth is no substitute for it.

\textbf{The eleven-term basis bought nothing.} Neither the two terms motivated
by the literature survey (exposure, signalling) nor the four motivated by
inspection improved on v0.3's three. The signalling term is worth singling out
because its sign was the interesting question: $+0.25$ (advertise a holding in a
half-suit our team is winning) scored $-0.300$ and $-0.25$ (conceal it) scored
$+0.056$, so if anything the concealment direction is favoured, but neither
interval comes close to excluding zero.

\subsection{The near-misses, and the one that resolved}
\label{sec:res-replication}

Two cells in Table~\ref{tab:nulls} had positive point estimates whose intervals
only just included zero, which is exactly the situation in which a screening
study manufactures a false positive if the best cell is promoted without
re-measurement. We reran both at 500 pairs on fresh seeds --- minimum detectable
effect $0.48$ rather than $0.77$ --- together with their combination:

\begin{center}
\small
\begin{tabular}{lrrl}
\toprule
candidate & screen (200 pairs) & retest (500 pairs) & retest 95\% CI \\
\midrule
blended value objective, $0.5$ & $+0.410$ & $-0.208$ & $[-0.534, +0.118]$ \\
$\sqrt{\cdot}$ counting, $\gamma = 0.70$ & $+0.480$ & $+0.212$ & $[-0.092, +0.516]$ \\
\textbf{both together} & --- & $\mathbf{+0.490}$ & $\mathbf{[+0.184, +0.796]}$ \\
\quad replication, fresh seeds & --- & $+0.068$ & $[-0.254, +0.390]$ \
\quad replication again & --- & $-0.066$ & $[-0.395, +0.263]$ \\
\bottomrule
\end{tabular}
\end{center}

The combination is the only cell in this entire study, out of more than
twenty-five, whose interval excluded zero in the improvement direction. On a
fresh set of 500 deals it returned $+0.068$, and on a second fresh set
$-0.066$. It was noise.

We report this at length because it is the single most useful methodological
observation the study produced. Twenty-five screening cells at a nominal 95\%
level will produce roughly one apparent winner under a complete null, and this
one arrived with a respectable-looking interval, a plausible story (two
independently-motivated changes that stack, exactly as v0.3's two winning terms
had) and a point estimate twice the sum of its parts --- which should have been
the tell. Promoting it would have shipped a change worth nothing and, worse,
would have anchored the next version's baseline on it.

\subsection{Belief-space search: unresolved, and why that is the finding}
\label{sec:res-lookahead}

Section~\ref{sec:lookahead} describes the one search design the variance
diagnosis does not rule out, and predicts, in Proposition~\ref{prop:greedy},
that a version of it without the quota coupling would collapse to greedy. Both
halves were measured. Table~\ref{tab:lookahead} is the whole record, and the
honest summary of it is that we do not know whether this search helps.

% BEGIN lookahead-table (generated by scripts4/lookahead_table.py --write)
\begin{table}[t]
\centering
\small
\begin{tabular}{lrrlr}
\toprule
cell & pairs & set diff & 95\% CI & MDE \\
\midrule
depth 3, $w=0.25$ & 200 & $\mathbf{+0.570}$ & $[+0.125, +1.015]$ & 0.77 \\
depth 2, $w=0.25$ & 200 & $+0.345$ & $[-0.083, +0.773]$ & 0.77 \\
depth 3, $w=0.60$ & 200 & $-0.080$ & $[-0.583, +0.423]$ & 0.77 \\
depth 3, \emph{no quota coupling} & 200 & $+0.070$ & $[-0.369, +0.509]$ & 0.77 \\
\midrule
depth 3, $w=0.25$ --- retest & 500 & $-0.044$ & $[-0.321, +0.233]$ & 0.48 \\
\midrule
\quad --- retest again & 500 & $\mathbf{+0.386}$ & $[+0.086, +0.686]$ & 0.48 \\
\midrule
\emph{no coupling} --- retest & 500 & $+0.112$ & $[-0.164, +0.388]$ & 0.48 \\
\midrule
depth 3, $w=0.25$ --- decisive A & 700 & $\mathbf{+0.307}$ & $[+0.063, +0.551]$ & 0.41 \\
\midrule
\quad --- decisive B & 700 & $+0.001$ & $[-0.240, +0.243]$ & 0.41 \\
\bottomrule
\end{tabular}
\caption{Belief-space lookahead against the v0.4 champion, which is the identical policy with the bonus disabled. MDE is the smallest effect that cell could have resolved at 80\% power given the measured per-pair standard deviation of 3.869 sets, so a null row says the run did not detect an effect of that size, not that there is none.}
\label{tab:lookahead}
\end{table}
% END lookahead-table

\paragraph{Three runs of one configuration disagree.} The screening cell, at 200
pairs, returned $+0.570$ and excluded zero. Following the replication policy of
Section~\ref{sec:res-replication} we reran the identical configuration twice
more at 500 pairs, changing nothing but the seed streams. The two retests
returned $-0.044$, 95\% CI $[-0.321, +0.233]$, and $+0.386$, 95\% CI
$[+0.086, +0.686]$. The second of those excludes zero. The first does not, and
its interval excludes the value the second reports.

That is not a result about the lookahead. It is a result about the experiment.
Three estimates of one quantity, differing only in seeds, are mutually
inconsistent: Cochran's $Q = 7.03$ on 2 degrees of freedom, $p = 0.030$, with
$I^2 = 71.5\%$ --- so roughly seven-tenths of the spread between these runs is
\emph{not} the sampling variation their own intervals model.

\paragraph{The within-run intervals are not the problem.} The per-pair standard
deviation recovered from each cell separately is $3.21$, $3.17$ and $3.42$ sets.
Those agree with each other and with the $3.869$ figure used throughout this
paper, so each interval correctly describes the deals that run actually played.
The natural inference is that the variance the intervals miss sits \emph{between}
runs; the DerSimonian--Laird estimate of such a component from these three cells
is $\tau = 0.268$ sets per deal-pair. \emph{We measured that inference directly
and it is wrong} --- see below --- but it is worth following through first,
because the two pooled answers it produces are the reason the question mattered
enough to measure.

Pooling all 1200 pairs on that assumption gives two answers, and only one of
them is usable:
\begin{center}
\begin{tabular}{lll}
\toprule
pooling & estimate & \\
\midrule
fixed effect   & $+0.226$ $[+0.041, +0.411]$ & excludes zero \\
random effects & $+0.278$ $[-0.083, +0.638]$ & includes zero \\
\bottomrule
\end{tabular}
\end{center}
A fixed-effect pool assumes the runs are three samples of one true value, which
is precisely what $Q$ appears to reject; the random-effects interval, which adds
$\tau^2$ to every cell, would then be the one reflecting where a fourth run might
land, and it includes zero. Read that way the effect is unresolved and could be
reported neither as working nor as a null --- which is where we left it, and
which is what made the size of $\tau$ worth establishing rather than assuming.

\paragraph{The three cells prompted a question about the harness, and the
answer is that the harness is fine.} The reading above --- that a duplicate-deal
cell has a variance component its own interval does not model --- would, if true,
make the minimum-detectable-effect table of Section~\ref{sec:eval} optimistic for
every single-cell result in this paper. That is too consequential to leave
inferred from three points, so we measured it.

The instrument is an \textbf{A/A}: one policy against a copy of itself, where the
true differential is exactly $0$. That required a change to the harness, and the
reason is worth recording, because it means this question could not previously
have been asked. Under the default seat-based seeding an A/A plays a
bit-identical game in both halves of a duplicate pair, so its differential is
$(a-b) + (b-a) = 0$ for every deal --- a point mass, not a small effect correctly
measured as zero. The harness could not observe its own null. Giving each side an
independent randomness stream (\texttt{independent\_seeds}, off by default so the
archive stays comparable) makes the null a real random variable.

Twenty-four blocks of 200 pairs, each with its own deal seeds and its own agent
seeds --- differing in exactly the way two separate runs differ:

\begin{center}
\begin{tabular}{lr}
\toprule
nominal 95\% intervals covering the truth & $23/24 = 95.8\%$ \ \ $[79.8, 99.3]$ \\
Cochran's $Q$ & $18.41$ on 23 df, $p = 0.735$ \\
$I^2$ & $0.0\%$ \\
$\tau$, between-run standard deviation & $\mathbf{0.0000}$ \\
grand mean of the 24 estimates & $-0.0215$ \ (truth: $0$) \\
per-pair standard deviation, 4800 pairs & $3.796$ \ (this paper uses $3.869$) \\
\bottomrule
\end{tabular}
\end{center}

Coverage is nominal. The observed spread of the block estimates, $0.2454$, is
\emph{smaller} than their own standard errors predict, $0.2701$; there is no
excess variance to attribute. A duplicate-deal interval means what it says.

\paragraph{So we retract the between-run variance.} An earlier version of this
section reported $\tau = 0.268$ sets per pair from the three cells above and
argued that it made every small result in this paper more fragile than it looks.
That was a Type I error. The $Q = 7.03$, $p = 0.030$ which produced it is exactly
the one-in-thirty event, and we promoted it to a structural claim about the
harness on the strength of three points --- which is the error this paper devotes
Section~\ref{sec:res-replication} and two of its four lessons to warning against,
committed inside the section doing the warning. The minimum-detectable-effect
table stands as written.

\paragraph{Which leaves the lookahead genuinely open, on better terms.} If the
cells are homogeneous --- and a harness that produces no excess heterogeneity
under the null is strong evidence that they are --- then fixed-effect pooling is
the valid analysis, and over all 1200 pairs the estimate is $+0.226$, 95\% CI
$[+0.041, +0.411]$.

We did not report that as a result. It excludes zero by $0.04$; this
configuration had already produced one screening cell that failed to replicate;
and reaching for the pooling that happens to resolve, immediately after the
pooling that happened not to, is the same error in the opposite direction. So we
fixed a size in advance and ran it: two further cells of 700 pairs each, on fresh
seeds, pre-registered as decisive.

They returned $+0.307$, 95\% CI $[+0.063, +0.551]$, and $+0.001$, 95\% CI
$[-0.240, +0.243]$.

\paragraph{The verdict, on the test we committed to.} Pooled, those two cells
give $+0.153$, 95\% CI $[-0.019, +0.324]$ --- which \emph{includes zero}. The
experiment designed in advance to settle this does not settle it, and we report
that rather than the alternative below, because a post-hoc pool does not get to
overrule a pre-registered one.

\paragraph{What the fuller picture says, and why we still do not claim it.}
Dropping the original screening cell --- it was \emph{selected} for resolving, so
including it in any pooled estimate is winner's curse --- and pooling the four
unselected cells over 2400 pairs gives $+0.153$, 95\% CI $[+0.022, +0.285]$. The
point estimate is identical to three decimal places across both analyses, which
is fairly strong evidence that a small real effect exists, of order a seventh of
a set per deal-pair. It is not evidence that we have demonstrated one.

Dropping that cell also disposes of the heterogeneity this section opened with:
$Q$ falls to $p = 0.063$ over the four unselected cells and $p = 0.081$ over the
two decisive ones, neither significant. Combined with the A/A measurement above,
the story is simply that one cell was selected for being large and the rest is
sampling noise --- no between-run variance, and nothing wrong with the harness.

\paragraph{And the sizing was wrong, by the error this section is about.} We
chose $n = 2600$ on the grounds that its minimum detectable effect of $0.21$
would separate $+0.23$ from zero. But $+0.23$ was the pool \emph{including} the
selected cell --- the inflated estimate. Against the unbiased $+0.153$, a
correctly powered test needs roughly $5000$ pairs. We committed the winner's
curse in the power calculation for the experiment whose purpose was to correct
for the winner's curse, which is the third time in this study the same error has
appeared one level up from where it was last caught.

The implementation ships disabled, which remains the right default for an effect
that is probably real and is not demonstrated.

\paragraph{What survives regardless.} Proposition~\ref{prop:greedy} does not
depend on any of this. A possession-chain search over a belief whose transition
edits only the asked card's row is exactly greedy at every depth, so the obvious
version of this idea is a no-op by construction rather than by measurement. That
is a structural statement about why search keeps failing in this game, it
removes a whole branch of the design space rather than a point in it, and it is
the part of Section~\ref{sec:lookahead} we would keep. Its empirical shadow is the
no-coupling cell, and the shadow is worth stating carefully, because the
proposition entails greedy equivalence \emph{at a node} and the step from there
to a null set differential is a modelling inference rather than a theorem.
Setting \texttt{couple=False} also does not restore the proposition's hypotheses
outright --- the transition still zeroes the asked row and decrements the
target's count. What is demonstrated is narrower and still to the point: with
the coupling off, the arm's chosen branch was the highest-probability branch in
$53{,}273$ of $53{,}273$ multi-branch nodes, against $49$ non-greedy choices with
it on; and its set differential is a null at both 200 and 500 pairs ($+0.070$
and $+0.112$, both intervals comfortably containing zero), which is consistent
with that. We would also withdraw a remark from an earlier draft that the
coupling ``changed no decision at all'': over $1{,}110$ champion-trajectory
positions the coupled and uncoupled arms disagree on $0.54\%$ of decisions, so
observing no disagreement in a sample of $59$ is the modal outcome rather than
evidence of equality.

The implementation ships disabled (\texttt{w\_lookahead} defaults to $0$), which
is the right default for an effect we cannot demonstrate.

% =====================================================================
\subsection{Is the opponent model just fitting its own kind?}
\label{sec:res-robust}

The model is a hypothesis about how opponents choose, and the opponents it was
measured against in Table~\ref{tab:gamma} are copies of the same policy. Two
tests separate a real effect from self-play overfitting.

\begin{table}[t]
\centering
\small
\begin{tabular}{llrrl}
\toprule
opponent & agent & pairs & set diff / pair & 95\% CI \\
\midrule
\multirow{2}{*}{\texttt{memory} (off-distribution)}
  & $\gamma = 0.30$ & 150 & $+4.727$ & $[+4.119, +5.335]$ \\
  & $\gamma = 0$    & 150 & $+3.793$ & $[+3.247, +4.340]$ \\
\midrule
\multirow{2}{*}{anti-model exploiter}
  & $\gamma = 0.30$ & 200 & $+6.570$ & $[+6.113, +7.027]$ \\
  & $\gamma = 0$    & 200 & $+5.735$ & $[+5.237, +6.233]$ \\
\bottomrule
\end{tabular}
\caption{The opponent model against opponents it was not shaped by.
\texttt{memory} acts only on facts it can prove and has no probabilistic
component at all. The anti-model exploiter is the v0.4 policy with the depth
term \emph{inverted} ($w_{\mathrm{suit}} = -0.35$), so it deliberately asks in
half-suits it is shallow in --- the exact behaviour the model assumes does not
happen.}
\label{tab:robust}
\end{table}

Against \texttt{memory}, a policy with no probabilistic component whatsoever,
the model is still worth about $+0.93$ sets per pair. Against an opponent built
specifically to violate its assumption it is still worth about $+0.84$. The
second number deserves care: the exploiter is a weak policy in absolute terms
(inverting the depth term costs it a great deal), so what the comparison shows
is that the model is not \emph{harmed} by an opponent designed to mislead it,
not that it cannot be exploited by a stronger adversary. A genuinely coordinated
best-responder is future work.

\subsection{Does it survive the rule variants?}
\label{sec:res-rules}

The engine supports several rule options, and a result that only holds under one
configuration is a weaker result. Table~\ref{tab:rules} repeats the headline
comparison under three of them.

\begin{table}[t]
\centering
\small
\begin{tabular}{llrrl}
\toprule
rule & comparison & pairs & set diff & 95\% CI \\
\midrule
declare at any time & v0.4 vs v0.3 champion & 200 & $+2.175$ & $[+1.647, +2.703]$ \\
declare at any time & $\gamma$ vs no $\gamma$ & 200 & $+1.935$ & $[+1.450, +2.420]$ \\
48-card deck & v0.4 vs v0.3 champion & 200 & $+1.080$ & $[+0.557, +1.603]$ \\
wrong split $\to$ opponents & v0.4 vs v0.3 champion & 200 & $+1.690$ & $[+1.191, +2.189]$ \\
\bottomrule
\end{tabular}
\caption{The result holds under every rule variant tested, including the
declare-at-any-time rule as the game is usually written down. Note the 48-card
row: with eight half-suits instead of nine, the margin roughly halves, which is
what one would expect if the advantage is inferential --- a smaller deck leaves
less to infer.}
\label{tab:rules}
\end{table}

The declare-at-any-time rows required an engine change. v0.3's game loop only
ever asked the player on move, so although the rule was accepted by the
configuration it could not be exercised by any policy --- and it matters most
precisely where it is unreachable, for a seat that has run out of cards and
therefore never gets a turn of its own while a teammate still has some. v0.4
polls off-turn seats for a claim they can make \emph{by deduction alone}, which
costs a belief update rather than a posterior per seat, and offers it identically
to both generations so that the comparison stays a comparison of one thing.

\subsection{Learning the ask objective}
\label{sec:res-learn}

v0.3's stated top priority for this version was to learn the ask objective
rather than hand-weight it. We did, by common-random-numbers paired rollout
regression: 874 positions, 83{,}168 games played to the end, all pairs of
candidate asks within a position regressed on the difference of their feature
vectors, with standard errors clustered by position (the naive standard errors
are on average $2.00\times$ too small, since pairs from one position share
candidates, worlds and rollout seeds).

The machinery works: common random numbers cut the variance of a pairwise
comparison to $0.64$ of independent worlds, bringing the noise-to-signal ratio
from v0.3's $2.4$ to $1.19$; a permutation test over position-level sign flips
gives $p = 0.003$ against a null $R^2$ of $0.005$; and a torch MLP on the same
paired target improves validation $R^2$ from $0.234$ to $0.249$, so the
objective's \emph{shape} is not the limitation either.

The learned weights lose. Played against the incumbent over 120 duplicate
deal-pairs, the best of three variants scored $-2.183$ sets per pair
$[-2.746, -1.621]$.

The diagnosis is specific, and it is a calibration failure of the target rather
than of the fit. Three of the eleven coefficients have answers already
established by \emph{play}: v0.3's 600-pair sweeps put \texttt{turn} at $0.60$
and \texttt{scarce} at $0.20$ and measured the pair as worth $+1.41$ sets per
deal-pair. The regression returns \texttt{turn} $= 0.157$ with a clustered
standard error of $0.058$ --- so an effect the size play has already measured
would sit ten standard errors from zero, and the dataset does not see it. More
starkly, position-centred rollout value rises by only $+0.101$ sets across the
entire range of $\Prob(\text{success})$, and an unpinned fit puts that
coefficient at $0.063$.

The likeliest mechanism is the continuation policy. A rollout has to be finished
by a policy that can attach to a determinized mid-game position, and the
belief tracker cannot: it is anchored on the initial deal and refuses. That
leaves a public-information heuristic, which throws away most of the value of a
marginal card, so a card won by a good ask is largely squandered before the game
ends and the ask stops mattering to the final differential. The binding
constraint is not precision --- pairing already fixed that --- it is that the
continuation is too weak for the difference between two asks to survive to the
end of the game.

This is the same shape as v0.3's search failure, one level up. There, the
problem was that the evaluation target could not rank asks; here, a much more
careful attempt to build a better target hits the same wall, and locates it in
the rollout policy rather than in the statistics.

\subsection{Perpetual states, measured}
\label{sec:res-perpetual}

Two hundred games per column, with ground truth used for analysis only.

\begin{table}[t]
\centering
\small
\begin{tabular}{lrrr}
\toprule
& normal & no progress rule & signalling on \\
\midrule
actions per game & 104.4 & 104.4 & 104.8 \\
games hitting the 4{,}000-action cap & \textbf{0} & \textbf{0} & \textbf{0} \\
games in which a position repeated & \textbf{200} & \textbf{200} & \textbf{200} \\
most repeats of a single position & 8 & 8 & 8 \\
games reaching a \emph{dead} position & \textbf{200} & \textbf{200} & \textbf{200} \\
share of plies fully dead & 4.24\% & 4.24\% & 4.56\% \\
actions played after going dead & 4.4 & 4.4 & 4.8 \\
games where a team held an unplaceable set & 110 & 110 & 116 \\
dead half-suits at peak & 3.13 & 3.13 & 3.15 \\
nulls per game & 0.275 & 0.275 & \textbf{0.220} \\
\bottomrule
\end{tabular}
\caption{Repetition and dead positions in 200 games of champion self-play.}
\label{tab:perpetual}
\end{table}

\paragraph{Repetition is universal, not exotic.} A position repeated in
\emph{every one} of 200 games, and some position recurred eight times in a
single game. v0.3 established cyclicity by search, as a possibility; it is in
fact a routine feature of ordinary play.

\paragraph{Every game dies before it ends.} All 200 games reached a position in
which no ask anywhere on the board could succeed. It happens late --- about 4.4
actions from the end, roughly 4.2\% of plies --- but it is universal. The last
phase of a game of Fish is not a contest at all; it is a bookkeeping exercise in
which the sets are already decided and only the declarations remain.

\paragraph{Termination is not the progress rule's doing.} With the agents'
stalling rule fully disabled, not one game in 200 hit the action cap, and the
mean game length was identical to four significant figures. So although the
rules genuinely permit non-termination, and although positions genuinely repeat,
in practice a strong policy walks out of every cycle on its own. The convention
is insurance, not the mechanism. Consistent with this, widening the stall window
from 80 to 200 actions produced 367 exact ties in 400 duplicate deals: the rule
almost never fires.

\paragraph{The deadlock is real, and it is expensive.} In 110 of 200 games
(55\%) a team reached a position where it could prove it held an entire
half-suit and still could not place the split. Those half-suits are nulled
\textbf{17.5\%} of the time against \textbf{2.8\%} for half-suits that never get
stuck --- a factor of $6.3$ --- and they account for \textbf{27\% of all nulls}
despite being only 11\% of half-suits.

\paragraph{The free-signalling escape works, and does not pay.} Turning on the
protocol of Section~\ref{sec:perpetual} cuts nulls from $0.275$ to $0.220$ per
game, a 20\% reduction, which confirms the mechanism does what the theory says.
It buys no wins: against the identical policy with signalling off, over 500
duplicate deals, $+0.002$ sets per pair $[-0.086, +0.090]$; restricted to
strictly dead positions, $+0.012$ $[-0.000, +0.024]$. Those intervals are far
tighter than the usual $\pm 0.48$ because the change is a near-no-op --- it alters
so few decisions that most pairs come out identical.

Two things are going on, and both are worth stating. The strictly-free version
of the protocol applies only when the \emph{whole board} is dead, which the seat
can prove on about 0.3\% of plies --- too rare to matter. The relaxed version
fires more often but is no longer free: it spends a turn to send one bit, and
the tempo it concedes cancels the sets it saves. So the escape is real theory
with a null in practice, and the honest summary is that Fish's stalemate is a
genuine structural feature of the game that a strong engine is already handling
adequately by other means.

\subsection{Game statistics, and a statistic that stops working}
\label{sec:res-analytics}

Homogeneous tables of each champion, 100 games each, on the measures v0.3 used
to characterise skill:

\begin{center}
\small
\begin{tabular}{lrr}
\toprule
& v0.3 champion & v0.4 champion \\
\midrule
cards gained per possession & 1.419 & 1.295 \\
possessions gaining nothing & 44.5\% & 48.8\% \\
failed-ask share & 40.9\% & 43.2\% \\
asks per game & 100.3 & 96.5 \\
ask accuracy, thirds of the game & 54.1 / 59.7 / 64.2\% & 52.2 / 58.0 / 60.8\% \\
null rate (natural half-suits) & 4.4\% & 3.1\% \\
claim delay, median / p90 (actions) & 0 / 58 & 0 / 49 \\
first-move advantage & $+0.23$ $[-0.27, +0.73]$ & $+0.32$ $[-0.22, +0.86]$ \\
\bottomrule
\end{tabular}
\end{center}

Three of these reproduce v0.3 and one contradicts it.

Reproduced: there is still no measurable first-move or seat advantage, both
intervals comfortably containing zero with the starting seat rotated across
deals; the synthetic ``eights and jokers'' half-suit still behaves like any
other (mean resolution rank $4.08$ against $3.99$, null rate $3.0\%$ against
$3.1\%$); and ask accuracy still rises from the opening into the later game,
so the opening remains the least-informed phase.

Contradicted, and worth stating carefully: v0.3 named cards-per-possession
``the single best summary statistic of skill'', on the evidence that it doubled
from the \texttt{memory} tier to the \texttt{probabilistic} tier. The stronger
engine here scores \emph{lower} on it. The mechanism is visible in the other
rows: v0.4 nulls fewer sets ($3.1\%$ against $4.4\%$) and claims sooner (p90 $49$
against $58$), so half-suits resolve faster and the game contains fewer, harder
asks --- $96.5$ per game against $100.3$. Cards-per-possession compares two
tables on how easy their asks were, and how easy the asks are depends on how
quickly the claiming policy is removing the resolved material. Within a tier
that shares a claim policy it tracks skill; across engines whose claiming
differs, it does not. A statistic that inverts when the thing it is measuring
improves is not a measure of skill, and we would retract the v0.3 claim rather
than qualify it.

% =====================================================================
\section{The engine as an instrument}
\label{sec:engine}

A strong policy that can only be measured in aggregate is a weak research
instrument, and an unusable one. The three things a chess engine gives a user
all exist here; they just mean something different in a game whose hidden state
is the entire problem.

\paragraph{Evaluation, in sets.} The banked differential plus, for every
unresolved half-suit, the learned probability that our team scores it minus the
probability that theirs does \eqref{eq:hsvalue}. So $+1.4$ reads as ``expect to
finish 1.4 sets ahead of the current score''. Because the underlying model is
fitted to outcomes and checked for calibration --- held-out log-loss $0.727$
against $0.821$ for the class prior, mean bias $-0.006$, reliable in every
decile --- the number means what it says. This is the one component that is a
good \emph{predictor} and a disastrous \emph{objective}
(Section~\ref{sec:res-null}), which is itself worth knowing: an evaluation
function and an objective function are not the same object.

\paragraph{A ranked move list with attribution.} Every legal ask, with the
probability it lands, the evaluation on each branch, and a per-term breakdown of
the policy's score. The breakdown matters for a research instrument rather than
a toy: when a human disagrees with the engine, the disagreement can be
attributed to a named term instead of left as an oracle's verdict.

\paragraph{A principal variation.} In chess the PV is the expected line. The
Fish analogue is the \emph{possession}: the chain of asks the engine would make
if each landed, since a turn in Fish is a run of successful asks. It is computed
by pinning each card as it is won and re-ranking, so it is the engine's actual
plan rather than a static list.

Two things a chess engine never has to show, and this one must. The
\textbf{card table} is the posterior over who holds every unresolved card ---
where the strength of this version actually lives, so it is displayed rather
than hidden. The \textbf{declaration table} gives, per live half-suit, the
probability our team holds all six, the most likely split, and the exact
probability that split is right, which is the quantity the claim decision turns
on.

The information boundary survives the user interface. Every panel is built from
the seat's own \texttt{Observation}; the posterior shown is an inference from the
public record, not a look behind it. The only exception is the reveal after the
last card is public anyway.

\paragraph{Mixed tables.} The same server seats any mixture of people and
engines at one table, which is what makes the engine usable as an instrument on
human play rather than only on self-play. The design point worth recording is
that this is not a cosmetic choice: teams in Fish are the alternating seats
$\{0,2,4\}$ against $\{1,3,5\}$, so a request as ordinary as ``three of us
against three bots'' is a constraint on \emph{which} seats the people take, not
on how many. A table therefore takes a seat \emph{arrangement} rather than a bot
count --- one that fills a whole side before crossing, and one that interleaves
--- and the 3-versus-3 configuration that a group of friends actually wants is
the former at $n=3$. Engine seats move on a clock rather than instantly, because
a decision costs about $4$\,ms and an unpaced possession appears to a human as a
single discontinuous jump. Pacing alone is not enough: a longer gap is
only a longer gap unless there is something to read in it, so the table
states each move in words together with what it \emph{proved}. That
second part is free, and is the one place where the no-bluff rule pays
the observer directly: an ask is a truthful statement about the asker's
hand as well as a question about the target's, so a failed ask
establishes that the target does not hold the card, that the asker does
not either, and that the asker holds some other card of that half-suit.
The default gap is $12$\,s and is adjustable in play. Freezing the table
turns the same control into a step-through --- one engine move per
request, still frozen afterwards --- which is what makes an engine's
whole possession readable, since a possession is a run of successful
asks and it is the run, not any single ask, that carries the plan. Each seat is served only its own
\texttt{Observation}, so the information boundary of
Appendix~\ref{app:leak} is enforced by the transport, not merely respected by
the client.

Analysis costs 130--270\,ms, which is interactive, and the same objects drive
the terminal client, the browser client and the batch experiments, so what a
user sees is what the experiments measured.

% =====================================================================
\section{Perpetual states: does Fish have a stalemate?}
\label{sec:perpetual}

Chess has stalemate and threefold repetition. The natural question is whether
Literature has an analogue --- a position the game cannot get out of --- and the
answer comes in three layers, of which only the third is interesting.

\subsection{Layer 1: the state graph is cyclic}

Nothing in the rules forces progress. Two opponents can trade a card back and
forth and return to a bit-identical position, so positions can repeat. This is a
property of the rules, not of any implementation. It follows immediately that
\emph{termination is a property of the players}: any claim that ``every game of
Fish ends'' describes a convention rather than the game. Our policies carry an
explicit progress rule and our harness a cap, and Section~\ref{sec:res-perpetual}
measures how much of ``games end'' those conventions are actually doing.

\subsection{Layer 2: under perfect information it does not matter}

Cycles exist, but the exact solver of Section~\ref{sec:exact} shows that no
\emph{optimal} line needs one: with all cards visible, optimal play always
terminates. Cyclicity on its own is therefore not a strategic phenomenon. The
reason is the closed form \eqref{eq:closedform} --- footholds only shrink, so the
mover simply collects everything reachable and the position marches forward.

\subsection{Layer 3: dead positions}

This layer has no chess counterpart, and it is the one worth having.

\begin{proposition}[Dead half-suits]
\label{prop:dead}
Call a live half-suit \emph{dead} when all six of its cards sit with one team.
Then no ask in that half-suit can ever succeed.
\end{proposition}

\begin{proof}
A legal ask in half-suit $H$ requires the asker to hold a card of $H$. If one
team holds all six, no member of the other team holds any card of $H$, so no
opponent of the holding team may ask in it at all. And a member of the holding
team may only ask an \emph{opponent}, who by hypothesis holds no card of $H$, so
such an ask fails. No card of $H$ therefore ever changes hands again.
\end{proof}

Call a position \emph{dead} when every live half-suit is dead. By
Proposition~\ref{prop:dead} the material is then frozen for the remainder of the
game: every legal ask is guaranteed to fail, the turn merely circulates, and the
only state changes left are declarations. This is as close to a stalemate as
Fish gets --- and, unlike a chess stalemate, it is not a draw. The sets are
already decided; they have simply not been claimed.

\subsection{The apparent deadlock, and why it is not one}

A dead position looks like a genuine trap for a team that holds a half-suit but
cannot place the split among its own members. No card will move again, so
--- the argument runs --- no new information can arrive, so the team must either
gamble on a declaration or stall forever.

That argument is wrong, and the way it is wrong is a strategy.

Information does not only arrive when cards move. Under the no-bluff rule a
player may not ask for a card they hold, so \textbf{asking for a card proves
publicly that you do not hold it}. In a dead position that ask is guaranteed to
fail, which means it costs no material whatsoever: there is no card to lose, and
the turn it concedes is worthless because no opponent can win a card either. The
ask is a \emph{free one-bit broadcast to your partners}.

A team stuck in a dead position can therefore talk its way out. Each member asks
for cards of the contested half-suit that they do not hold; after a few such
asks the split is common knowledge within the team and the declaration is safe.
Literature's only communication channel is normally costly and read by the
opposition; in a dead position it becomes free, and the opposition can do
nothing with what it hears, because the half-suit is already beyond their reach.

\begin{remark}[A definition that made the phenomenon invisible]
Our first implementation defined ``our team provably holds this half-suit'' as
``every card's holder is pinned and all six are teammates''. That makes ``we hold
it but cannot place the split'' a contradiction in terms, since pinning the
holders \emph{is} placing the split, and the condition never fired once in 1,047
measured plies. The correct condition is weaker and is the whole point: every
\emph{possible} holder of every card is a teammate, which does not require
knowing which teammate. Under that definition the situation occurs on 2.6\% of
plies.
\end{remark}

% =====================================================================
\section{What this says about playing Fish}
\label{sec:strategy}

The results have readings that do not require an engine, and we state them
because a study of a game people play ought to say something to the people
playing it.

\paragraph{Watch where opponents ask, not just what they ask for.} This is the
finding of the paper, and it is actionable at the table. A player who asks in a
half-suit is probably deep in it, in rough proportion to how many of its cards
they hold, and that single inference is worth more than every other improvement
we made combined. Concretely: when an opponent asks in hearts, do not only note
the hard fact that they hold at least one heart --- shift weight toward their
holding several. The engine's version of this is one parameter, and it is worth
about $1.9$ sets per duplicate deal-pair against a strong opponent that ignores
it.

\paragraph{An exact rule for the perfect-information endgame.} Once every
remaining card is public, the game is solved in closed form: the team on move
scores every live half-suit in which it holds at least one card, and the
opponents score the rest. Footholds only shrink and a well-informed asker never
misses. The practical corollary is that in a fully-known endgame the only thing
that matters is \emph{how many live half-suits your team has a foothold in when
it is your turn} --- not how many cards, not which cards.

\paragraph{Do not agonise over when to declare.} Reproducing v0.3's finding on a
much better posterior: raising the claim bar from $0.97$ to $0.99$ changes
essentially no decisions. The reason is structural rather than statistical. If
your team genuinely holds all six cards of a half-suit, no opponent holds any
card of it, so no opponent can legally ask in it --- and a claim by a team that
does not hold all six gives the set to you. Waiting is close to free. Wait until
you know the split.

\paragraph{Two tie-breakers really do hold up.} When two asks look about equally
likely to land, prefer the one that, if it misses, hands the turn to the
opponent with fewer cards, and prefer to fight in half-suits your team is
already winning. v0.3 measured that pair at $+1.41$ sets per deal-pair; with an
inference layer good enough to make the primary criterion much sharper, they are
still worth $+1.08$. Better card-reading does not make positional judgement
redundant.

\paragraph{The endgame is bookkeeping, so do the bookkeeping.} Every one of 200
measured games reached a position where no ask anywhere could succeed: the sets
were decided and only the declarations were left. Two consequences for a human.
First, the last few turns are not a contest, so spend your attention earlier.
Second --- and this is the one people get wrong --- if your team holds a half-suit
outright, \emph{no opponent can legally ask in it}, so it cannot be taken from
you. There is never a reason to rush that declaration. In our games, half-suits
that reached "ours but we cannot place the split" were nulled 17.5\% of the
time against 2.8\% otherwise, and accounted for 27\% of all nulls. Almost every
one of those is a self-inflicted wound.

\paragraph{You can talk, and most players do not realise it.} You may not ask
for a card you hold. So asking for a card \emph{proves to the table} that you do
not have it --- and when your team already owns the half-suit, the ask cannot
lose you anything, because no opponent holds a card of it to give. In that
situation an ask is a free, legal message to your partners telling them exactly
where a card is not. Use it when you are stuck on a split. (Our engine gains
about 20\% fewer nulls doing this; it does not gain games, because the turn it
spends is worth about as much as the sets it saves. For a human, who is far more
likely to misplace a split than an engine is, the trade is probably better.)

\paragraph{And a warning about the obvious idea.} Scoring asks by ``how much
does this improve my team's chance of winning this half-suit'' --- which sounds
like the right objective, is measured in the right units, and rests on a
well-calibrated model --- is a disaster: $-7.4$ sets per deal-pair against simply
asking for the card most likely to be there. In this game, the probability that
the ask lands is not one consideration among several. It is the objective, and
everything else is a tie-break.

% =====================================================================
\section{Failed and abandoned experiments}
\label{sec:failures}

Recorded deliberately, because in this project the negative results have
repeatedly been the informative ones.

\paragraph{Exact inference, on its own.} The headline effort of this version ---
replacing a documented-biased sampler with provably correct inference --- bought
nothing in play (Section~\ref{sec:res-exact}) and was measurably worse at
predicting reality (Section~\ref{sec:res-posterior}). We report it as a failure
of the \emph{hypothesis}, not of the implementation: the implementation is
validated against exhaustive enumeration in four separate ways. What failed was
the assumption that the uniform posterior is the right target.

\paragraph{Rejection sampling for the disjunctive clauses.} Exactly correct, and
unaffordable: measured acceptance mean $0.30$, median $0.18$, below $0.05$ on
16\% of positions. Abandoned in favour of importance sampling with closed-form
densities.

\paragraph{Folding the clauses into the exact dynamic program.} Sound, and
tractable in the median case (refined state space $5{,}760$), but with a tail
that makes it unusable as a default: p90 $230{,}000$ and p99
$2.7 \times 10^{7}$. Retained as an option, not as the path.

\paragraph{A na\"ive batch rewrite of the sampler.} Vectorising the draw across
the batch gave only $1.4\times$ until the dictionary materialisation of the
draws was removed as well; the correct fix was to carry the whole batch as an
integer matrix end to end, after which the same change gave $5.1\times$. Worth
recording because the first measurement would have justified abandoning a
change that was in fact the right one.

\paragraph{An under-converged fit that looked plausible.} The first half-suit
value model was fitted by plain gradient descent on unstandardised features and
reported a held-out log-loss \emph{worse than the class prior} while producing
entirely reasonable-looking coefficients. Feature scales spanning two orders of
magnitude were the cause. This is the reason the fitting routine now reports
against two baselines (the class prior and a single-feature model) rather than
against nothing.

\paragraph{A false positive, caught.} Of more than twenty-five screening cells,
exactly one produced an interval excluding zero in the improvement direction
($+0.490$, $[+0.184, +0.796]$, 500 pairs). It failed to replicate twice on fresh
seeds ($+0.068$ and $-0.066$). We record it as a failed experiment rather than
burying it in a methodology note, because it is the shape of the mistake this
project is most likely to make next.

\paragraph{Style that responds to the match.} Every objective term in v0.4 is a
function of one candidate ask and the posterior. None knows the score, and none
knows that the last four moves were the same two players passing one card back
and forth. Two ideas were built to close that gap, each ablatable to zero, and
both are refuted --- not as nulls, but with a dose-response, which is the more
informative outcome because the gradient confirms the mechanism rather than
merely the sign. All cells are 200 duplicate deal-pairs against the champion.

\begin{center}
\small
\begin{tabular}{lrl}
\toprule
change & set diff & 95\% CI \\
\midrule
break the duel, $w_{\mathrm{retake}} = 0.30$ & $-0.340$ & $[-0.665, -0.015]$ \\
break the duel, $w_{\mathrm{retake}} = 1.00$ & $\mathbf{-1.015}$ & $[-1.562, -0.468]$ \\
\midrule
tie-breakers up when behind, $w_{\mathrm{behind}} = 0.50$ & $-0.185$ & $[-0.628, +0.258]$ \\
tie-breakers up when behind, $w_{\mathrm{behind}} = 1.00$ & $\mathbf{-0.890}$ & $[-1.384, -0.396]$ \\
tie-breakers down when behind, $w_{\mathrm{behind}} = -0.50$ & $-0.075$ & $[-0.552, +0.402]$ \\
\bottomrule
\end{tabular}
\end{center}

\emph{Breaking the duel} penalises taking back a card this seat just lost to that
same opponent. Section~\ref{sec:game} notes that the state graph is cyclic
precisely because two opponents can trade a card back and forth; the question
here was whether a player should decline to. They should not, and three times the
penalty costs three times the sets. This is the scheduling argument of
Proposition~\ref{prop:greedy} arriving from a third direction: a card you watched
an opponent take is a card you \emph{know} they hold, so taking it back is a
certain ask and a certain ask keeps the turn. Deferring it risks the turn to
preserve an option that was not going anywhere, and the information concealed by
not asking is worth less than the possession thrown away. It is worth recording
that this is a manoeuvre strong players describe using, which the engine could
not previously express at all --- the measurement is the point, not the outcome.

\emph{Score-adaptive tie-breakers} scale the weights by how far behind the team
is, on the reasoning that a losing side should buy information and a winning side
should buy certainty. Scaling them \emph{up} when behind loses monotonically.
Scaling them \emph{down} is indistinguishable from the champion. The asymmetry is
the finding: the loss surface is steep on the side of trusting the tie-breakers
more and flat on the side of trusting them less, which is what one would expect
if v0.3's hand-tuned constants already sit at the point where the tie-breakers
have given everything they have --- and is a third piece of evidence, after the
inverted-U sweeps and the substitution result of Section~\ref{sec:res-ladder},
that they are tie-breakers rather than an objective.

Neither idea is retried in another form. What both share is that they spend
$P(\text{success})$ to buy something else, and this paper's most reproducible
finding is that nothing else is worth buying with it.

\paragraph{Belief-space possession-chain search.} Built to test the one design
the variance diagnosis leaves open, and abandoned \emph{unresolved} rather than
as a null. Three runs of the identical configuration returned $+0.570$ (200
pairs), $-0.044$ and $+0.386$ (500 pairs each); the first and third exclude zero
and the second excludes the third's estimate. Under random-effects pooling,
which is the only valid pooling given $Q = 7.03$ on 2 df, the answer is
$+0.278$, 95\% CI $[-0.083, +0.638]$. We could have reported this as a success or
as a failure by choosing a run. Over the four unselected cells --- 2400 pairs, the
selected screening cell excluded --- it is $+0.153$, 95\% CI $[+0.022, +0.285]$,
while the two cells pre-registered as decisive pool to $[-0.019, +0.324]$ and
include zero. A small effect, probably real, not demonstrated; and the decisive
run was under-powered because its size was chosen against an estimate the
selected cell had inflated. Section~\ref{sec:res-lookahead} has the arithmetic.
The implementation ships disabled.

\paragraph{Inherited from v0.3, and not revisited.} Determinized search (PIMC
and Information-Set MCTS) lost to its own prior; a value network trained on
belief features and applied inside determinized worlds lost 34--6 to a
train/inference distribution mismatch; quiescence extension made search worse;
and an expected-value claim model predicted an optimal threshold near $0.70$
that direct measurement falsified. None of these were re-attempted here, because
v0.3's diagnosis --- that the binding constraint is the objective, not the depth
of search --- was the premise of this version and is not contradicted by
anything we measured.

% =====================================================================
\section{Limitations}
\label{sec:limitations}

\paragraph{The opponent model is a model.} Everything else in the inference
stack is deduction, validated against enumeration. Equation~\eqref{eq:oppmodel}
is a hypothesis about how people choose, fitted with one parameter. Against
opponents whose ask policy is unrelated to their depth in a half-suit it could
in principle hurt, and against an opponent who deliberately inverts the
heuristic it is exploitable in the ordinary sense that any model is. We measure
both cases rather than assert robustness, and $\gamma$ is a single number that
can be set to zero.

\paragraph{Depth is evaluated on the initial deal.} The choice model is really
about what a player held \emph{when they asked}. Cards move, so the two diverge
late in a game. The exact quantity would need a replay of the public history per
candidate world, turning an $O(1)$ update into $O(|h|)$.

\paragraph{The uniform-posterior result is about \emph{this} game and
\emph{these} opponents.} We show that exact inference under an uninformative-
choice assumption is worse than a biased sampler that accidentally encodes a
choice model. That is a statement about which approximation is closer to the
truth, not a general licence to prefer biased inference.

\paragraph{Exact solving is still bounded.} Absolute ground truth reaches
positions with one, and now two, unresolved half-suits. Midgame positions remain
out of reach, which is exactly where v0.3 localised the remaining gap.

\paragraph{No equilibrium claim.} We compute no equilibrium and no
exploitability bound. The appropriate solution concept for a game with pre-play
agreement and no in-play communication is TMECor
\citep{celli2018computational}, and nothing here approximates it. Every strength
statement is relative to a named opponent or absolute against an exactly solved
subgame.

\paragraph{A rules discrepancy, stated.} The printed rules allow a set to be
declared at any point, including during an opponent's turn. v0.3's game loop
only ever asked the player on move, so the rule was accepted by the engine but
unreachable by any policy --- which matters most for a seat that has run out of
cards and therefore never gets a turn of its own while a teammate still has
some. v0.4 adds off-turn declaration, offered identically to both generations
so that the comparison stays a comparison of one thing. The primary evaluation
is nevertheless run under v0.3's own configuration, so that the numbers are
comparable with the previous version, with the variant reported separately.

% =====================================================================
\section{Conclusions}
\label{sec:conclusions}

We set out to remove a defect the previous version had documented in its own
code: a belief sampler that was not uniform over the deals consistent with the
public record. We removed it, exactly. Hidden-state inference in Literature
reduces to counting assignments in a degree-constrained bipartite system, and
although that is $\#P$-complete in general, the instances the game produces are
tiny in the dimension that matters --- roughly four distinct candidate sets even
when twenty-six cards are unlocated --- so a dynamic program returns the exact
partition function, exact marginals, exact uniform sampling and exact joint
probabilities. The disjunctive constraints that the program cannot absorb are
handled by importance sampling with closed-form proposal densities, verified
against exhaustive enumeration for support, density and marginals, with a
control demonstrating that dropping the weights visibly breaks the estimate.

And it changed nothing. Against the previous champion, at matched seeds over 250
duplicate deal-pairs, provably correct inference moved the set differential by
$-0.008$ per pair.

The instructive part is why, and that we could find out. Because the hidden
hands exist in simulation, a posterior can be scored on the question it is
actually trying to answer, and that measurement says the exactly-uniform
posterior is \emph{worse} at predicting reality than the biased sampler it
replaced. Proposition~\ref{prop:uniform} is a correct theorem about a false
hypothesis: it assumes players' choices carry no information beyond legality.
The old sampler's bias, an artefact of how it seeded the disjunctive
constraints, happened to encode a crude model of how players choose. Removing
the bias removed the model with it.

Replacing the accident with one line of modelling --- a player asks in a
half-suit in proportion to how many cards of it they hold --- recovers the loss
and goes past it, on all three instruments we have: posterior likelihood against
the true hands, agreement with an exact solver, and duplicate-deal play. The
last of these puts it at $+1.9$ sets per deal-pair against the previous
champion, which is larger than the entire objective improvement that separated
that champion from its own baseline.

Three lessons generalise beyond this game.

\textbf{Correct is not the same as good.} An inference procedure is correct
relative to a model. When a cheaper, wronger procedure beats an exact one, the
exact procedure's \emph{assumptions} are the thing that is failing, and the
productive response is to find which assumption and relax it --- not to add
compute to the correct-but-misspecified thing.

\textbf{Measure the intermediate quantity, not only the outcome.} Every belief
metric in the previous version was indirect, and a change that was provably
correct and empirically useless would have been very hard to diagnose from win
rates alone. Scoring the posterior directly against ground truth turned a
puzzling null result into an identified cause in one experiment. Where a
simulator gives access to the hidden state, refusing to use it for
\emph{evaluation} is a false economy; the discipline that matters is not letting
it reach the acting policy.

\textbf{Know what your experiment could not have found.} The measured per-pair
standard deviation of the set differential here is $3.869$ sets, which means the
previous version's 600-pair cells could not resolve effects below about $0.44$
sets per pair. Several of its reported nulls were therefore uninformative rather
than negative. We report the minimum detectable effect alongside every interval
for that reason.

\textbf{Replicate the winner, not the losers.} A screening study over
twenty-five cells will hand you an apparent winner under a complete null, and it
will arrive with a plausible mechanism attached --- ours came with a story about
two independently-motivated changes stacking, which is exactly what the previous
version had genuinely found for a different pair of terms. The cheap defence is
not a multiplicity correction, it is one more run on fresh seeds. Ours cost forty
minutes and saved a wrong champion.

\subsection{What we would do next}

\begin{itemize}
\item \textbf{A proper likelihood, not a proxy.} The choice model uses depth in
the asked half-suit on the initial deal. The right object is the acting
policy's own probability of the observed action given the hand, evaluated at the
time of the ask. That is a fixpoint --- the model of the opponent is the policy
--- and it is expensive, but it is the principled version of the thing that
already works.
\item \textbf{A coordinated best-responder.} The robustness tests here show the
model is not harmed by an opponent built to violate its assumption, which is
weaker than showing it cannot be exploited. Training a genuine best response,
and reporting the value it extracts, is the only route to an absolute claim
about how good the champion is.
\item \textbf{Cross-play.} Once anything in this engine is \emph{learned} from
self-play rather than configured, cross-play between independently seeded runs
becomes necessary rather than optional, for the reason Hanabi established: an
agent that only wins alongside copies of itself has not solved the game.
\item \textbf{Belief-space search, past the one point we tested.}
Section~\ref{sec:lookahead} built the design this section previously proposed
--- a lookahead over the posterior rather than over sampled determinizations ---
and Section~\ref{sec:res-lookahead} leaves it open. Settling it needs pairs
rather than a cleverer analysis --- the harness's intervals are measured to be
honest, so 1200 pairs simply are not enough to separate $+0.23$ from $0$ --- and
the question is narrow either way: it covers possession-chain value, a single
sweep of quota coupling, beam $4$, depth $3$, entering additively. Proposition~\ref{prop:greedy} says the
version without the coupling could never have worked, which removes a whole
branch of the space rather than one point in it. What is still open is a leaf
evaluator other than banked cards, a converged rather than single-sweep
coupling, and a horizon that does not stop at the end of one possession.
\end{itemize}

% =====================================================================
\appendix

\section{Reproducing this work}
\label{app:repro}

The engine, agents, experiments and this paper are in one repository. v0.3's
code is untouched, so its results remain reproducible alongside these; v0.4
lives in \texttt{fish4/}.

\begin{verbatim}
py -m pytest tests  -q     # v0.3: rules, fuzz, leakage proofs, belief
                           # soundness, statistics, exact solver
py -m pytest tests4 -q     # v0.4: exact counting vs brute force, sampler
                           # support/density/marginals, batch equivalence,
                           # information-boundary invariance, exact solver v2

py scripts4/duel.py jobs/j6_ladder.json 4      # the strength ladder
py scripts4/duel.py jobs/j11_final.json 6      # headline + ablations
py scripts4/summarise_duels.py                 # every duel, with its MDE

py scripts4/posterior_accuracy.py 14 3   # posteriors scored on the true hands
py scripts4/exact_bench4.py              # agreement with provably optimal play
py scripts4/fit_hsvalue.py 200 3         # fit the half-suit value model
py scripts4/analytics4.py 120 3          # strategy statistics
py scripts4/check_tex.py                 # structural check on this document
\end{verbatim}

Every duel appends to \texttt{results/v04\_duels.jsonl} recording both random
seed streams, the full agent specifications, the rule configuration and the
result. \texttt{summarise\_duels.py} prints, alongside every interval, the
smallest effect that run could have detected --- which is the column that stops
a null being read as a negative.

\section{The champion}
\label{app:champion}

\begin{verbatim}
from fish4.registry4 import make_agent
agent = make_agent(("fishbot4", {"opponent_gamma": 0.35}))
\end{verbatim}

Everything else is at its default, and the defaults are deliberately v0.3's own
ask objective ($w_{\mathrm{suit}} = 0.06$, $w_{\mathrm{turn}} = 0.60$,
$w_{\mathrm{scarce}} = 0.20$). The improvement reported here comes from
inference, not from re-tuning the objective on top of it --- and
Section~\ref{sec:res-null} is the record of trying to re-tune it and failing.

\section{Information-boundary tests}
\label{app:leak}

The property asserted is that every quantity v0.4 computes is bit-identical when
the hidden cards of the other five seats are permuted into a different layout
equally consistent with the public record. Candidate permutations are drawn from
the acting seat's own belief sampler --- uniform permutation almost never lands
on a consistent world, because the public record is very constraining --- and
each candidate is then validated by an \emph{independent} replay of the whole
public history, so the test does not assume the sampler is correct. Checked
quantities: posterior marginals under four combinations of the opponent-model
parameters, all eleven ask-objective terms, and the end-to-end policy action.

There is a second boundary, and it is worth recording that the first set of
tests did not protect it. The permutation tests establish that the
\emph{policy} never reads a card it should not; they say nothing about what the
\emph{server} puts on the wire. An earlier revision of the multiplayer table
included the deal seed in the per-seat view, as a debugging convenience. Every
hand is a deterministic function of that integer, so one number defeated the
entire property the permutation tests were built to guarantee --- and it was
readable without joining the table, since a table code alone fetches a view.
The lesson generalises past this bug: a leak test that only exercises the
component you were thinking about will keep passing while the leak moves to the
component you were not. The transport now has its own test, which walks the
serialised view recursively and fails on any key or value that could reconstruct
the deal, for a seated player and for an anonymous visitor alike.

% =====================================================================
\begin{thebibliography}{99}

\bibitem[Celli and Gatti(2018)]{celli2018computational}
Andrea Celli and Nicola Gatti.
\newblock Computational results for extensive-form adversarial team games.
\newblock In \emph{AAAI Conference on Artificial Intelligence}, 2018.

\bibitem[Long et al.(2010)]{long2010understanding}
Jeffrey Richard Long, Nathan R. Sturtevant, Michael Buro, and Timothy Furtak.
\newblock Understanding the success of perfect information Monte Carlo sampling
in game playing.
\newblock In \emph{AAAI Conference on Artificial Intelligence}, pages 134--140,
2010.

\bibitem[Bard et al.(2020)]{bard2020hanabi}
Nolan Bard, Jakob N. Foerster, Sarath Chandar, Neil Burch, Marc Lanctot,
H.~Francis Song, et al.
\newblock The Hanabi challenge: a new frontier for AI research.
\newblock \emph{Artificial Intelligence}, 280:103216, 2020.

\bibitem[Bez\'akov\'a et al.(2012)]{bezakova2012negative}
Ivona Bez\'akov\'a, Alistair Sinclair, Daniel \v{S}tefankovi\v{c}, and Eric
Vigoda.
\newblock Negative examples for sequential importance sampling of binary
contingency tables.
\newblock \emph{Algorithmica}, 64(4):606--620, 2012.

\bibitem[Chen et al.(2005)]{chen2005sequential}
Yuguo Chen, Persi Diaconis, Susan P. Holmes, and Jun S. Liu.
\newblock Sequential Monte Carlo methods for statistical analysis of tables.
\newblock \emph{Journal of the American Statistical Association},
100(469):109--120, 2005.

\bibitem[Cowling et al.(2012)]{cowling2012ismcts}
Peter I. Cowling, Edward J. Powley, and Daniel Whitehouse.
\newblock Information set Monte Carlo tree search.
\newblock \emph{IEEE Transactions on Computational Intelligence and AI in
Games}, 4(2):120--143, 2012.

\bibitem[Frank and Basin(1998)]{frank1998search}
Ian Frank and David Basin.
\newblock Search in games with incomplete information: a case study using Bridge
card play.
\newblock \emph{Artificial Intelligence}, 100(1--2):87--123, 1998.

\bibitem[Goadrich et al.(2026)]{goadrich2026valet}
Mark Goadrich, Aiden Morenville, and \'Eric Piette.
\newblock Valet: a standardized testbed of traditional imperfect-information
card games.
\newblock arXiv:2603.03252, 2026.

\bibitem[Hu et al.(2020)]{hu2020otherplay}
Hengyuan Hu, Adam Lerer, Alex Peysakhovich, and Jakob N. Foerster.
\newblock ``Other-play'' for zero-shot coordination.
\newblock In \emph{International Conference on Machine Learning}, 2020.

\bibitem[Jerrum et al.(2004)]{jerrum2004polynomial}
Mark Jerrum, Alistair Sinclair, and Eric Vigoda.
\newblock A polynomial-time approximation algorithm for the permanent of a
matrix with nonnegative entries.
\newblock \emph{Journal of the ACM}, 51(4):671--697, 2004.

\bibitem[R\'egin(1996)]{regin1996generalized}
Jean-Charles R\'egin.
\newblock Generalized arc consistency for global cardinality constraint.
\newblock In \emph{AAAI Conference on Artificial Intelligence}, pages
209--215, 1996.

\bibitem[Rebstock et al.(2019)]{rebstock2019policy}
Douglas Rebstock, Christopher Solinas, Michael Buro, and Nathan R. Sturtevant.
\newblock Policy based inference in trick-taking card games.
\newblock In \emph{IEEE Conference on Games}, 2019.

\bibitem[Solinas et al.(2019)]{solinas2019improving}
Christopher Solinas, Douglas Rebstock, and Michael Buro.
\newblock Improving search with supervised learning in trick-based card games.
\newblock In \emph{AAAI Conference on Artificial Intelligence}, volume 33,
pages 1158--1165, 2019.

\bibitem[Solinas et al.(2023)]{solinas2023history}
Christopher Solinas, Douglas Rebstock, Nathan R. Sturtevant, and Michael Buro.
\newblock History filtering in imperfect information games: algorithms and
complexity.
\newblock arXiv:2311.14651, 2023.

\bibitem[von Stengel and Koller(1997)]{vonstengel1997team}
Bernhard von Stengel and Daphne Koller.
\newblock Team-maxmin equilibria.
\newblock \emph{Games and Economic Behavior}, 21(1--2):309--321, 1997.

\bibitem[Waudby-Smith and Ramdas(2024)]{waudbysmith2024estimating}
Ian Waudby-Smith and Aaditya Ramdas.
\newblock Estimating means of bounded random variables by betting.
\newblock \emph{Journal of the Royal Statistical Society Series B},
86(1):1--27, 2024.

\bibitem[Valiant(1979)]{valiant1979complexity}
Leslie G. Valiant.
\newblock The complexity of computing the permanent.
\newblock \emph{Theoretical Computer Science}, 8(2):189--201, 1979.

\end{thebibliography}

\end{document}
